\documentclass[11pt,a4paper]{article}

\usepackage[T1]{fontenc}
\usepackage[utf8]{inputenc}
\usepackage[margin=2.5cm]{geometry}
\usepackage{amsmath,amssymb,amsthm,mathrsfs}
\usepackage{mathtools}
\usepackage{bm}
\usepackage{enumitem}
\usepackage{booktabs}
\usepackage{xcolor}
\usepackage{hyperref}
\usepackage{cleveref}

\hypersetup{colorlinks=true,linkcolor=blue!60!black,citecolor=blue!60!black,urlcolor=blue!60!black}

\newcommand{\E}{\mathbb{E}}
\newcommand{\Var}{\mathbb{V}}
\newcommand{\dd}{\mathrm{d}}
\newcommand{\bbar}[1]{\bar{#1}}
\renewcommand{\ss}{\star}

\title{A Smets--Wouters like model}
\author{Technical Note}
\date{}

\begin{document}
\maketitle
\tableofcontents
\newpage

%% ============================================================
\section{Overview}
%% ============================================================

This note derives the equilibrium conditions of a medium-scale New Keynesian
DSGE model in the tradition of \textsc{Smets and Wouters}~(2003, 2007). The
model extends the canonical framework with the following features:
\begin{itemize}[nosep]
  \item \textbf{Kimball aggregation} for both goods and labour markets (variable demand elasticity);
  \item \textbf{External habit formation} with a stochastic habit parameter;
  \item \textbf{Non-separable preferences} between consumption and labour;
  \item \textbf{Variable capacity utilisation} with endogenous depreciation;
  \item \textbf{Investment adjustment costs} (Christiano--Eichenbaum--Evans type);
  \item \textbf{Fiscal instruments}: consumption tax, labour income tax, income tax, employer labour tax;
  \item \textbf{Calvo pricing and wage setting} with partial indexation;
  \item A rich set of \textbf{17 structural shocks}.
\end{itemize}
All variables are expressed in \emph{detrended} form, i.e.\ divided by the
stochastic labour-augmenting technology level~$A_t^T$. The gross growth rate of
technology is $g_t \equiv A_t^T / A_{t-1}^T$, with steady-state value
$1 + \bbar{g}$.


%% ============================================================
\section{Notation and conventions}
%% ============================================================

\begin{center}
\begin{tabular}{lll}
\toprule
\textbf{Symbol} & \textbf{Dynare name} & \textbf{Description} \\
\midrule
$C_t$         & \texttt{Consumption}          & Detrended consumption \\
$L_t$         & \texttt{HouseholdLabourSupply} & Household labour supply \\
$\lambda_t$   & \texttt{HouseholdLagrangeMultiplier} & Marginal utility of wealth \\
$I_t$         & \texttt{Investment}            & Detrended investment \\
$\bbar K_t$   & \texttt{CapitalStock}          & Physical (installed) capital stock \\
$\delta_t$    & \texttt{DepreciationRate}       & Depreciation rate \\
$\delta'_t$   & \texttt{dDepreciationRate}      & Marginal depreciation rate \\
$S_t$         & \texttt{InvestmentCost}         & Investment adjustment cost \\
$S'_t$        & \texttt{dInvestmentCost}        & Marginal investment adj.\ cost \\
$z_t$         & \texttt{CapacityUtilizationFactor} & Capacity utilisation rate \\
$r^k_t$       & \texttt{CapitalReturnRate}      & Real rental rate of capital \\
$w^h_t$       & \texttt{HouseholdRealWage}      & Household real wage \\
$Q_t$         & \texttt{TobinQ}                 & Tobin's $Q$ \\
\midrule
$w_t$         & \texttt{RealGrossWage}          & Gross real wage \\
$L^d_t$       & \texttt{LabourDemand}           & Labour demand \\
$K_t$         & \texttt{CapitalDemand}          & Effective capital (services) \\
$mc_t$        & \texttt{RealMarginalCost}       & Real marginal cost \\
$\breve p_t$       & \texttt{OptimalRelativePrice}   & Optimal relative price \\
$Z_{1,t},Z_{2,t},Z_{3,t}$ & \texttt{Z1,Z2,Z3}  & Price recursion variables \\
$Y_t$         & \texttt{GDP}                    & Detrended output \\
$\Lambda^f_t$ & \texttt{FinalGoodLagrangeMultiplier} & Kimball price aggregator \\
$\pi^{np}_t$  & \texttt{NonOptimizingFirmsPriceGrowth} & Non-optimiser price growth \\
$\tilde p_t$  & \texttt{AveragedRelativePrices} & Average relative price \\
$\Delta^p_t$  & \texttt{PriceDistorsion}        & Price dispersion \\
$\nabla^p_t$  & \texttt{NablaPrice}             & Price recursion variable \\
$\pi_t$       & \texttt{InflationFactor}         & Gross inflation \\
\midrule
$\omega_t$    & \texttt{RealWageGrowthFactor}   & Real wage growth \\
$\tilde w_t$  & \texttt{AveragedRelativeWages}  & Average relative wage \\
$\breve w_t$       & \texttt{OptimalRelativeRealWage} & Optimal relative wage \\
$H_{1,t},H_{2,t},H_{3,t}$ & \texttt{H1,H2,H3}  & Wage recursion variables \\
$\Lambda^s_t$ & \texttt{EmploymentAgencyLagrangeMultiplier} & Kimball wage aggregator \\
$L^s_t$       & \texttt{EmploymentAgencyLabourSupply} & Employment agency supply \\
$\omega^{nw}_t$ & \texttt{NonOptimizingUnionsWageGrowth} & Non-optimiser wage growth \\
$\Delta^w_t$  & \texttt{WageDistorsion}         & Wage dispersion \\
$\nabla^w_t$  & \texttt{NablaWage}              & Wage recursion variable \\
\midrule
$R_t$         & \texttt{NominalInterestFactor}  & Gross nominal interest rate \\
$G_t$         & \texttt{PublicSpending}          & Government spending \\
$\mathcal{G}_t$ & \texttt{OutputGap}            & Output gap \\
\bottomrule
\end{tabular}
\end{center}

\medskip

\noindent\textbf{Exogenous processes.}
\begin{center}
\begin{tabular}{lll}
\toprule
\textbf{Symbol} & \textbf{Dynare name} & \textbf{Description} \\
\midrule
$g_t$ & \texttt{ProductionEfficiencyGrowth} & Technology growth rate (gross) \\
$A_t$ & \texttt{ProductionEfficiencyCycle} & Stationary technology level \\
$\tau^C_t$ & \texttt{ConsumptionTax} & Consumption tax \\
$\tau^W_t$ & \texttt{LabourIncomeTax} & Labour income tax (employee) \\
$\tau^R_t$ & \texttt{IncomeTax} & Income tax \\
$\tau^L_t$ & \texttt{LabourTax} & Labour tax (employer) \\
$\varepsilon^B_t$ & \texttt{RiskPremium} & Risk premium shock \\
$\varepsilon^L_t$ & \texttt{LabourSupplyShock} & Labour supply shock \\
$p^I_t$ & \texttt{InvestmentRelativePrice} & Investment relative price \\
$\varepsilon^I_t$ & \texttt{InvestmentEfficiencyShock} & Investment efficiency \\
$\nu^p_t$ & \texttt{PriceCostPushShock} & Price cost-push shock \\
$\nu^w_t$ & \texttt{WageCostPushShock} & Wage cost-push shock \\
$\bbar\pi_t$ & \texttt{InflationTarget} & Inflation target \\
$\varepsilon^R_t$ & \texttt{TaylorShock} & Monetary policy shock \\
$\varepsilon^g_t$ & \texttt{PublicSpendingShare} & Public spending share \\
$\eta^s_t$ & \texttt{HabitShock} & Habit shock \\
$\varsigma_t$ & \texttt{PreferenceShock} & Preference shock \\
\bottomrule
\end{tabular}
\end{center}

\noindent\textbf{Key structural parameters.}
\begin{center}
\begin{tabular}{lll}
\toprule
\textbf{Symbol} & \textbf{Name} & \textbf{Description} \\
\midrule
$\beta$ & discount factor & subjective discount factor \\
$\sigma_c$ & risk aversion & inverse of the intertemporal elasticity of substitution \\
$\sigma_l$ & labour curvature & inverse of the Frisch elasticity \\
$\eta$ & habit & degree of external habit formation \\
$\psi$ & investment cost & investment adjustment cost parameter \\
$\alpha$ & capital share & capital elasticity in production \\
$\xi_p$ & price Calvo & probability of \emph{not} re-optimising price \\
$\xi_w$ & wage Calvo & probability of \emph{not} re-optimising wage \\
$\gamma_p$ & price indexation & degree of price indexation to past inflation \\
$\gamma_w$ & wage indexation & degree of wage indexation to current inflation \\
$\theta_f$ & price elasticity & demand elasticity parameter (goods) \\
$\psi_f$ & Kimball curvature (goods) & Kimball demand curvature (goods) \\
$\theta_s$ & wage elasticity & demand elasticity parameter (labour) \\
$\psi_s$ & Kimball curvature (labour) & Kimball demand curvature (labour) \\
$\mu_f$ & price markup & steady-state goods markup \\
$\mu_s$ & wage markup & steady-state wage markup \\
\bottomrule
\end{tabular}
\end{center}


%% ============================================================
\section{Households}
\label{sec:households}
%% ============================================================

\subsection{Preferences}

A representative household maximises expected lifetime utility:
\begin{equation}\label{eq:utility}
  \E_0 \sum_{t=0}^{\infty} \beta^t \, u(C_t, L_t; \bm{\xi}_t),
\end{equation}
where the period utility function takes the \textbf{non-separable} form:
\begin{equation}\label{eq:period_utility}
  u(C_t, L_t; \bm{\xi}_t)
  = \frac{1}{1-\sigma_c}
    \Bigl(C_t - \eta\,\eta^s_t\,\frac{C_{t-1}}{g_t} + \ln\varsigma_t\Bigr)^{1-\sigma_c}
    \exp\!\Biggl(\frac{(\sigma_c - 1)\,\varepsilon^L_t}{1+\sigma_l}
    \Bigl(\frac{L_t}{\bbar L}\Bigr)^{1+\sigma_l}\Biggr).
\end{equation}

This specification embeds several key features:
\begin{itemize}[nosep]
  \item \textbf{External habit formation}: the effective consumption
    $\tilde C_t = C_t - \eta\,\eta^s_t \, C_{t-1}/g_t$ depends on lagged
    \emph{aggregate} consumption, scaled by the stochastic habit parameter
    $\eta^s_t$ and detrended by the technology growth factor~$g_t$.
  \item \textbf{Preference shock}: $\ln\varsigma_t$ is an additive preference
    shifter.
  \item \textbf{Non-separability}: when $\sigma_c \neq 1$, the marginal
    utility of consumption depends on labour supply through the exponential
    term. The labour supply shock $\varepsilon^L_t$ and the term
    $(L_t/\bbar L)^{1+\sigma_l}$ enter multiplicatively, generating a wealth
    effect on labour supply.
\end{itemize}

\subsection{Budget constraint}

In each period, the household chooses consumption~$C_t$,
labour supply~$L_t$, investment~$I_t$, physical capital~$\bbar K_t$, capacity
utilisation~$z_t$, and bond holdings. The (detrended, real) flow budget
constraint is:
\begin{equation}\label{eq:budget}
  \tau^C_t\, C_t + p^I_t\,I_t
  + \frac{B_t}{R_t\,P_t}
  = \tau^R_t\bigl(\tau^W_t\, w^h_t\, L_t
  + r^k_t\, z_t\, \bbar K_{t-1}/g_t\bigr)
  + \frac{B_{t-1}}{P_t\,g_t}
  + \Pi_t + T_t,
\end{equation}
where:
\begin{itemize}[nosep]
  \item $\tau^C_t$ is the consumption tax;
  \item $p^I_t$ is the relative price of investment goods;
  \item $\tau^R_t$ is the income tax (applies to both labour and capital income);
  \item $\tau^W_t$ is a labour-specific income tax (employee side);
  \item $w^h_t$ is the household real wage (after the wage markup);
  \item $r^k_t$ is the rental rate on effective capital;
  \item $z_t$ is the capacity utilisation rate;
  \item $\Pi_t$ are profits from firms; $T_t$ are lump-sum transfers.
\end{itemize}

\noindent The household also faces a \textbf{capital accumulation} equation:
\begin{equation}\label{eq:capital_accum}
  \bbar K_t = (1 - \delta_t)\,\frac{\bbar K_{t-1}}{g_t}
  + \varepsilon^I_t \,(1 - S_t)\, I_t\,,
\end{equation}
where $\delta_t = \delta(z_t)$ is the endogenous depreciation rate and
$S_t = S(I_t, I_{t-1}, g_t)$ captures investment adjustment costs.

\subsection{Investment adjustment cost function}

The adjustment cost function depends on the growth rate of investment:
\begin{equation}\label{eq:adj_cost}
  S_t \equiv S\!\left(\frac{I_t}{I_{t-1}}\,g_t\right)
  = \frac{\psi}{2}\,(1+\bbar g)^2
    \left(\frac{I_t}{I_{t-1}}\,\frac{g_t}{1+\bbar g} - 1\right)^{\!2},
\end{equation}
with derivative
\begin{equation}\label{eq:adj_cost_deriv}
  S'_t \equiv \frac{\partial S}{\partial (I_t g_t / I_{t-1})}
  = \psi\,(1+\bbar g)
    \left(\frac{I_t}{I_{t-1}}\,\frac{g_t}{1+\bbar g} - 1\right).
\end{equation}
At steady state, $I_t/I_{t-1} = 1$ and $g_t = 1+\bbar g$, so
$S = S' = 0$: there are no adjustment costs along the balanced growth path.

\subsection{Endogenous depreciation}

The depreciation rate is an increasing, convex function of capacity
utilisation~$z_t$:
\begin{equation}\label{eq:depreciation}
  \delta(z_t) = \bbar\delta \, \exp\!\left(
    \frac{\bbar\delta'}{\bbar\delta}\,(z_t - \bbar z)\right),
\end{equation}
where $\bbar\delta = \delta(\bbar z)$ is the steady-state depreciation rate and
$\bbar\delta' = \delta'(\bbar z)$ is the steady-state marginal depreciation.
The marginal depreciation rate is:
\begin{equation}\label{eq:depreciation_deriv}
  \delta'(z_t) = \bbar\delta' \, \exp\!\left(
    \frac{\bbar\delta'}{\bbar\delta}\,(z_t - \bbar z)\right).
\end{equation}

\subsection{First-order conditions}

The household maximises~\eqref{eq:utility} subject to the budget
constraint~\eqref{eq:budget} and capital accumulation~\eqref{eq:capital_accum}.
Let $\lambda_t$ denote the Lagrange multiplier on the budget constraint and
$Q_t\,\lambda_t$ the multiplier on the capital accumulation equation (so that
$Q_t$ is Tobin's~$Q$). The first-order conditions yield the following
equations, numbered as in the Dynare code.

\medskip

\paragraph{Equation~1: Consumption FOC.}
Differentiating with respect to $C_t$:
\begin{equation}\label{eq:eq1}\tag{Eq.~1}
  \boxed{
  \bigl(C_t - \eta\,\eta^s_t\,C_{t-1}/g_t + \ln\varsigma_t\bigr)^{-\sigma_c}
  \exp\!\Bigl(\tfrac{(\sigma_c-1)\varepsilon^L_t}{1+\sigma_l}
  (L_t/\bbar L)^{1+\sigma_l}\Bigr)
  = \tau^C_t\,\lambda_t.}
\end{equation}
The left-hand side is the marginal utility of consumption. The right-hand side
is the shadow cost of one extra unit of consumption, inclusive of the
consumption tax.

\paragraph{Equation~2: Euler equation.}
The first-order condition for bond holdings, combined with the risk premium
shock~$\varepsilon^B_t$, gives:
\begin{equation}\label{eq:eq2}\tag{Eq.~2}
  \boxed{
  \lambda_t\,\tau^R_t
  = \beta\,\varepsilon^B_t\,R_t\,g_{t+1}^{-\sigma_c}\,
    \frac{\lambda_{t+1}\,\tau^R_{t+1}}{\pi_{t+1}}.}
\end{equation}
This is the standard Euler equation: the marginal cost of saving (left-hand
side, inclusive of the income tax) equals the discounted marginal benefit. The
term $g_{t+1}^{-\sigma_c}$ arises from the detrending of the utility function.

\paragraph{Equation~3: Labour supply FOC.}
Differentiating with respect to $L_t$:
\begin{equation}\label{eq:eq3}\tag{Eq.~3}
  \boxed{
  \begin{aligned}
  &\bigl(C_t - \eta\,\eta^s_t\,C_{t-1}/g_t + \ln\varsigma_t\bigr)^{1-\sigma_c}
  \,\frac{\varepsilon^L_t}{\bbar L}
  \,\Bigl(\frac{L_t}{\bbar L}\Bigr)^{\!\sigma_l} \\
  &\qquad\times \exp\!\Bigl(\tfrac{(\sigma_c-1)\varepsilon^L_t}{1+\sigma_l}
  (L_t/\bbar L)^{1+\sigma_l}\Bigr)
  = \lambda_t\,\tau^R_t\,\tau^W_t\,w^h_t.
  \end{aligned}}
\end{equation}
The left-hand side is the marginal disutility of labour (with a sign
convention making it positive). The right-hand side is the marginal benefit of
working one more unit, after all taxes.

\paragraph{Equation~4: Investment FOC.}
Differentiating with respect to $I_t$, using the capital accumulation
constraint:
\begin{equation}\label{eq:eq4}\tag{Eq.~4}
  \boxed{
  \begin{aligned}
  &Q_t\!\left(1 - S_t - \frac{I_t}{I_{t-1}}\,g_t\,S'_t\right)\\
  &\quad + \beta\,\frac{\lambda_{t+1}}{\lambda_t}\,g_{t+1}^{-\sigma_c}\,
    Q_{t+1}\,\frac{\varepsilon^I_{t+1}}{\varepsilon^I_t}\,
    \left(\frac{I_{t+1}}{I_t}\,g_{t+1}\right)^{\!2}\!S'_{t+1}
  = \frac{p^I_t}{\varepsilon^I_t}.
  \end{aligned}}
\end{equation}
The left-hand side captures the shadow value of an extra unit of investment
(current marginal value plus the expected future benefit from relaxing
tomorrow's adjustment cost). The right-hand side is the effective price of
investment.

\paragraph{Equation~5: Depreciation rate.}
This is the definition~\eqref{eq:depreciation}:
\begin{equation}\label{eq:eq5}\tag{Eq.~5}
  \boxed{
  \delta_t = \bbar\delta\,
  \exp\!\Bigl(\frac{\bbar\delta'}{\bbar\delta}(z_t - \bbar z)\Bigr).}
\end{equation}

\paragraph{Equation~6: Marginal depreciation rate.}
This is the definition~\eqref{eq:depreciation_deriv}:
\begin{equation}\label{eq:eq6}\tag{Eq.~6}
  \boxed{
  \delta'_t = \bbar\delta'\,
  \exp\!\Bigl(\frac{\bbar\delta'}{\bbar\delta}(z_t - \bbar z)\Bigr).}
\end{equation}

\paragraph{Equation~7: Capacity utilisation FOC.}
Differentiating with respect to $z_t$:
\begin{equation}\label{eq:eq7}\tag{Eq.~7}
  \boxed{\tau^R_t\,r^k_t = Q_t\,\delta'_t.}
\end{equation}
The household equates the after-tax marginal revenue from utilising capital
one unit more intensively (left-hand side) to the marginal cost in terms of
faster depreciation (right-hand side).

\paragraph{Equation~8: Capital stock FOC (Tobin's~$Q$).}
Differentiating with respect to $\bbar K_t$:
\begin{equation}\label{eq:eq8}\tag{Eq.~8}
  \boxed{
  Q_t = \beta\,\frac{\lambda_{t+1}}{\lambda_t}\,g_{t+1}^{-\sigma_c}
  \Bigl[Q_{t+1}(1-\delta_{t+1})
  + \tau^R_{t+1}\,r^k_{t+1}\,z_{t+1}\Bigr].}
\end{equation}
Tobin's~$Q$ equals the discounted expected value of capital tomorrow: the
resale value $(1-\delta_{t+1})Q_{t+1}$ plus the after-tax rental income
$\tau^R_{t+1}\,r^k_{t+1}\,z_{t+1}$.

\paragraph{Equation~9: Capital accumulation.}
This restates~\eqref{eq:capital_accum}:
\begin{equation}\label{eq:eq9}\tag{Eq.~9}
  \boxed{
  \bbar K_t = (1-\delta_t)\,\frac{\bbar K_{t-1}}{g_t}
  + \varepsilon^I_t\,(1-S_t)\,I_t.}
\end{equation}

\paragraph{Equation~10: Investment cost.}
This restates~\eqref{eq:adj_cost}:
\begin{equation}\label{eq:eq10}\tag{Eq.~10}
  \boxed{
  S_t = \frac{\psi}{2}\,(1{+}\bbar g)^2
    \Bigl(\frac{I_t}{I_{t-1}}\,\frac{g_t}{1{+}\bbar g} - 1\Bigr)^{\!2}.}
\end{equation}

\paragraph{Equation~11: Marginal investment cost.}
This restates~\eqref{eq:adj_cost_deriv}:
\begin{equation}\label{eq:eq11}\tag{Eq.~11}
  \boxed{
  S'_t = \psi\,(1{+}\bbar g)
    \Bigl(\frac{I_t}{I_{t-1}}\,\frac{g_t}{1{+}\bbar g} - 1\Bigr).}
\end{equation}


%% ============================================================
\section{Intermediate goods and price setting}
\label{sec:prices}
%% ============================================================

\subsection{Production technology}

A continuum of intermediate goods firms $\iota \in [0,1]$ each produce a
differentiated variety using capital services $K_t(\iota)$ and effective
labour $A_t L^d_t(\iota)$:
\begin{equation}\label{eq:production}
  y_t(\iota) = K_t(\iota)^{\alpha}\,(A_t\,L^d_t(\iota))^{1-\alpha},
\end{equation}
where $\alpha\in(0,1)$ is the capital share and $A_t$ is the stationary
component of technology.

\subsection{Cost minimisation}

Each intermediate firm minimises its real cost $r^k_t\,K_t(\iota) +
\tau^L_t\,w_t\,L^d_t(\iota)$ subject to~\eqref{eq:production}, where
$\tau^L_t$ is the employer-side labour tax (payroll tax).

\paragraph{Equation~13: Factor price frontier.}
The ratio of factor demands satisfies:
\begin{equation}\label{eq:eq13}\tag{Eq.~13}
  \boxed{\frac{\tau^L_t\,w_t\,L^d_t}{r^k_t\,K_t} = \frac{1-\alpha}{\alpha}.}
\end{equation}

\paragraph{Equation~14: Real marginal cost.}
Substituting optimal factor demands back into the cost function:
\begin{equation}\label{eq:eq14}\tag{Eq.~14}
  \boxed{
  mc_t = A_t^{\,\alpha-1}\,
  \Bigl(\frac{r^k_t}{\alpha}\Bigr)^{\!\alpha}\,
  \Bigl(\frac{\tau^L_t\,w_t}{1-\alpha}\Bigr)^{\!1-\alpha}.}
\end{equation}
Note that the marginal cost is identical across firms (since all face the
same factor prices), even though each firm has a different level of output.

\subsection{Final good production: Kimball aggregation}

A representative final-good firm assembles intermediate varieties using a
\textbf{Kimball aggregator} (Kimball, 1995). Unlike the standard
Dixit--Stiglitz CES aggregator, the Kimball aggregator generates a
\emph{variable} elasticity of demand: firms that raise their price far
above the market level face a much steeper fall in demand than under CES.
This strategic complementarity in price setting is an important ingredient
for matching the empirical persistence of inflation.

\subsubsection{The Kimball aggregation constraint}

The aggregation technology is defined implicitly by:
\begin{equation}\label{eq:kimball_goods}
  \int_0^1 G\!\left(\frac{y_t(\iota)}{Y_t};\,\theta_f,\psi_f\right)\dd\iota = 1,
\end{equation}
where $y_t(\iota)$ is the output of intermediate firm~$\iota$ and $Y_t$ is
aggregate output. Following Dotsey and King~(2005) and
Adjemian, Darracq Pari\`es and Moyen~(2008), the function~$G$ takes the
explicit form:
\begin{equation}\label{eq:kimball_G}
  \boxed{
  G(x;\theta_f,\psi_f) = \frac{\theta_f}{\theta_f(1+\psi_f)-1}
  \Bigl[(1+\psi_f)\,x - \psi_f\Bigr]
    ^{\frac{\theta_f(1+\psi_f)-1}{\theta_f(1+\psi_f)}}
  - \frac{\theta_f}{\theta_f(1+\psi_f)-1} + 1.}
\end{equation}
One can verify that $G(1) = 1$ (normalisation) and $G'(1) = 1$.

The two parameters have the following roles:
\begin{itemize}[nosep]
  \item $\theta_f > 1$ controls the \textbf{elasticity of demand} at the
    symmetric equilibrium. The steady-state goods markup is
    $\mu_f = \theta_f/(\theta_f-1)$, i.e.\
    $\theta_f = \mu_f/(\mu_f - 1)$.
  \item $\psi_f$ controls the \textbf{curvature} of the demand function
    (how the elasticity varies with the relative price). The Kimball
    curvature parameter $\varepsilon_f$ is defined by
    $\psi_f = -\varepsilon_f/\theta_f$.
    When $\psi_f = 0$, the aggregator reduces to the standard CES
    (Dixit--Stiglitz) case.
\end{itemize}

\subsubsection{Demand for variety~$\iota$}

The final-good firm maximises profit
$Y_t - \int_0^1 p_t(\iota)\,y_t(\iota)\,\dd\iota$ subject
to~\eqref{eq:kimball_goods}, where $p_t(\iota) = P_t(\iota)/P_t$ is the
relative price of variety~$\iota$. The first-order condition with respect to
$y_t(\iota)$ is:
\begin{equation}\label{eq:kimball_foc}
  p_t(\iota) = \Lambda^f_t\,G'\!\left(\frac{y_t(\iota)}{Y_t}\right),
\end{equation}
where $\Lambda^f_t$ is the Lagrange multiplier on the aggregation
constraint. Differentiating~\eqref{eq:kimball_G}:
\[
  G'(x) = \Bigl[(1+\psi_f)\,x - \psi_f\Bigr]^{-\frac{1}{\theta_f(1+\psi_f)}}.
\]
Substituting into~\eqref{eq:kimball_foc} and inverting for $y_t(\iota)/Y_t$:
\begin{equation}\label{eq:kimball_demand}
  \frac{y_t(\iota)}{Y_t}
  = \frac{1}{1+\psi_f}\left[
    \left(\frac{p_t(\iota)}{\Lambda^f_t}\right)^{\!-\theta_f(1+\psi_f)}
    + \psi_f\right].
\end{equation}
This is the \textbf{Kimball demand function}. When $\psi_f = 0$, it
reduces to the familiar CES demand $y_t(\iota)/Y_t =
(p_t(\iota))^{-\theta_f}$ (with $\Lambda^f_t = 1$).

\subsubsection{Variable elasticity of demand}

The price elasticity of demand at variety~$\iota$ is:
\[
  \varepsilon_t(\iota)
  \equiv -\frac{\partial\ln(y_t(\iota)/Y_t)}{\partial\ln p_t(\iota)}
  = \frac{\theta_f(1+\psi_f)\,
    (p_t(\iota)/\Lambda^f_t)^{-\theta_f(1+\psi_f)}}
    {(p_t(\iota)/\Lambda^f_t)^{-\theta_f(1+\psi_f)} + \psi_f}.
\]
At the symmetric equilibrium ($p_t(\iota) = \Lambda^f_t = 1$), this
simplifies to $\varepsilon = \theta_f$, confirming that $\theta_f$ is
the elasticity at the steady state. When $\psi_f < 0$ (the empirically
relevant case), the elasticity \emph{increases} as a firm raises its
price above the market level, discouraging large price deviations and
generating strategic complementarity.

\subsubsection{Aggregate conditions from the Kimball aggregator}

To derive the aggregate price relationships, define $x(\iota) = y_t(\iota)/Y_t$
and substitute the demand function~\eqref{eq:kimball_demand} into the
aggregation constraint~\eqref{eq:kimball_goods}. Integrating over all
varieties requires tracking three cross-sectional moments of relative
prices. Under Calvo pricing (where a fraction~$\xi_p$ of firms have
``inherited'' prices and a fraction~$1-\xi_p$ have the optimal
price~$\breve p_t$), these moments evolve recursively.

\paragraph{Average relative price $\tilde p_t$.}
The first moment of relative prices is $\tilde p_t =
\int_0^1 p_t(\iota)\,\dd\iota$. Under Calvo:
\begin{equation}\label{eq:eq20}\tag{Eq.~20}
  \boxed{
  \tilde p_t = (1-\xi_p)\,\breve p_t
  + \xi_p\,\frac{\pi^{np}_t}{\pi_t}\,\tilde p_{t-1}.}
\end{equation}
Optimising firms contribute $\breve p_t$; non-optimising firms had relative
price $\tilde p_{t-1}$ on average last period, and their nominal price
grew by $\pi^{np}_t$ while the aggregate price grew by $\pi_t$, so their
relative price changes by $\pi^{np}_t/\pi_t$.

\paragraph{Kimball multiplier $\Lambda^f_t$.}
From the aggregation constraint, the Lagrange multiplier satisfies:
\begin{equation}\label{eq:eq16}\tag{Eq.~16}
  \boxed{
  (\Lambda^f_t)^{1-\theta_f(1+\psi_f)}
  = (1-\xi_p)\,(\breve p_t)^{1-\theta_f(1+\psi_f)}
  + \xi_p\,
    \left(\frac{\pi^{np}_t}{\pi_t}\right)^{\!1-\theta_f(1+\psi_f)}
    \!(\Lambda^f_{t-1})^{1-\theta_f(1+\psi_f)}.}
\end{equation}
This generalises the CES price index. When $\psi_f = 0$,
$\Lambda^f_t = 1$ always and the equation reduces to the standard
Dixit--Stiglitz price aggregation.

\paragraph{Aggregate price normalisation.}
The zero-profit condition of the final-good firm, combined with the
Kimball aggregation, yields:
\begin{equation}\label{eq:eq21}\tag{Eq.~21}
  \boxed{
  \frac{\psi_f}{1+\psi_f}\,\tilde p_t + \frac{1}{1+\psi_f}\,\Lambda^f_t = 1.}
\end{equation}
This links the average relative price and the Kimball multiplier so that
the aggregate price level is normalised to unity. At steady state,
$\tilde p = \Lambda^f = 1$.

\subsubsection{Price dispersion}

We now derive the price dispersion index, which measures the resource
cost of relative price heterogeneity under the Kimball aggregator.

\paragraph{Definition of $\Delta^p_t$.}
The Kimball demand~\eqref{eq:kimball_demand} implies that the output
of intermediate firm~$\iota$ relative to aggregate output is:
\[
  \frac{y_t(\iota)}{Y_t}
  = \frac{1}{1+\psi_f}
  \left[\left(\frac{p_t(\iota)/\Lambda^f_t}{}\right)^{\!-(1+\psi_f)\theta_f}
    + \psi_f\right].
\]
Integrating over all firms $\iota \in [0,1]$:
\begin{equation}\label{eq:deltap_def}
  \Delta^p_t
  \equiv \int_0^1 \frac{y_t(\iota)}{Y_t}\,\dd\iota
  = \frac{1}{1+\psi_f}
  \left[(\Lambda^f_t)^{\theta_f(1+\psi_f)}\!
    \int_0^1 p_t(\iota)^{-\theta_f(1+\psi_f)}\,\dd\iota
    + \psi_f\right],
\end{equation}
where the integral $\int_0^1 p_t(\iota)^{-\theta_f(1+\psi_f)}\,\dd\iota$
is a ``power mean'' of relative prices that we
denote~$\nabla^p_t$.

\paragraph{Recursive expression for $\nabla^p_t$.}
Under the Calvo mechanism, at date~$t$ the price $P_t(\iota)/P_t$ has
been set optimally $j$ periods earlier with probability
$(1-\xi_p)\,\xi_p^j$. Since then it has been indexed by the cumulative
factor $\Gamma_{t,t-j}/\pi_{t-j+1}\cdots\pi_t$. Thus:
\[
  \nabla^p_t
  = \int_0^1 p_t(\iota)^{-\theta_f(1+\psi_f)}\,\dd\iota
  = (1-\xi_p)\sum_{j=0}^{\infty}\xi_p^j
    \left(\frac{\Gamma_{t,t-j}\,\breve P_{t-j}}{P_t}\right)^{\!-\theta_f(1+\psi_f)}.
\]
This infinite sum can be split into the $j=0$ term (newly optimising
firms) and the remaining terms ($j\geq 1$, which correspond to
$\nabla^p_{t-1}$ evaluated at prices deflated by one more period of
indexation):
\begin{equation}\label{eq:eq23}\tag{Eq.~23}
  \boxed{
  \nabla^p_t = (1-\xi_p)\,(\breve p_t)^{-\theta_f(1+\psi_f)}
  + \xi_p\,\nabla^p_{t-1}\,
    \left(\frac{\pi^{np}_t}{\pi_t}\right)^{\!-\theta_f(1+\psi_f)}.}
\end{equation}
The ratio $\pi^{np}_t/\pi_t$ captures the erosion of each
non-optimising firm's relative price: its nominal price grows by
$\pi^{np}_t$ (indexation) while the aggregate price level grows by
$\pi_t$.

\paragraph{Equation~22: Price dispersion index $\Delta^p_t$.}
Substituting the definition of $\nabla^p_t$ back
into~\eqref{eq:deltap_def}:
\begin{equation}\label{eq:eq22}\tag{Eq.~22}
  \boxed{
  \Delta^p_t = \frac{\psi_f}{1+\psi_f}
  + \frac{1}{1+\psi_f}\,\nabla^p_t\,
    (\Lambda^f_t)^{\theta_f(1+\psi_f)}.}
\end{equation}
The first term $\psi_f/(1+\psi_f)$ is the contribution of the linear
part of the Kimball demand (the $+\psi_f$ term). The second term involves
the power mean of relative prices $\nabla^p_t$ scaled by the Kimball
multiplier~$\Lambda^f_t$.

In the CES case ($\psi_f = 0$), the expression simplifies to
$\Delta^p_t = \nabla^p_t = \int_0^1 p_t(\iota)^{-\theta_f}\,\dd\iota$,
the standard Yun~(1996) price dispersion index.

At steady state, $\breve p=1$, $\pi^{np}=\pi$, $\Lambda^f=1$, so
$\nabla^p = 1$ and $\Delta^p = 1$: there is no dispersion when all
firms charge the same price.

\subsection{Calvo pricing and optimal price setting}

In each period, a fraction $(1-\xi_p)$ of intermediate firms can
re-optimise their price. The remaining fraction~$\xi_p$ adjusts
mechanically according to an indexation rule.

\paragraph{Equation~12: Price indexation for non-optimising firms.}
Non-optimising firms set their price growth as a geometric average of
lagged inflation and the inflation target:
\begin{equation}\label{eq:eq12}\tag{Eq.~12}
  \boxed{
  \pi^{np}_t = \pi_{t-1}^{\gamma_p}\,\bbar\pi_t^{\,1-\gamma_p},}
\end{equation}
where $\gamma_p \in [0,1]$ is the degree of backward-looking indexation.

\subsubsection{The firm's profit and the Calvo lottery}

The nominal profit of a firm proposing price~$\mathcal{P}$ at date~$t$
is:
\[
  \Pi_t(\mathcal{P}) =
  \left(\frac{\mathcal{P}}{P_t} - \nu^p_t\,mc_t\right)
  \left[\left(\frac{\mathcal{P}/P_t}{\Theta_t}\right)^{\!-(1+\psi_f)\theta_f}
    + \psi_f\right]
  \frac{P_t\,Y_t}{1+\psi_f},
\]
where the term in square brackets is the Kimball
demand~\eqref{eq:kimball_demand} evaluated at $\mathcal{P}/P_t$,
$mc_t$ is the (common) real marginal cost, and $\nu^p_t$ is the
price cost-push shock ($\nu^p_t > 1$ raises effective marginal cost).

The cost-push shock~$\nu^p_t$ (on the cost side) is the model's
structural proxy for the Smets and Wouters ``price markup shock.''  It
enters only~$Z_1$ and therefore affects the Phillips curve without
distorting the resource constraint; we make this precise in
Section~\ref{sec:linearised_nkpc}.

The firm is subject to a lottery \`a la Calvo~(1983): at each date it
receives a signal indicating whether it can or cannot re-optimise its
price. With probability~$\xi_p$ the signal is unfavourable and the firm
simply indexes its price according to the \emph{ad hoc} rule:
\begin{equation}\label{eq:index_rule}
  P_t(\iota) = \Gamma_t\,P_{t-1}(\iota),
  \qquad\text{where}\quad
  \Gamma_t \equiv \pi_{t-1}^{\gamma_p}\,\bbar\pi^{\,1-\gamma_p}.
\end{equation}
More generally, if a firm has not re-optimised for $j$ consecutive
periods, its price has been multiplied by
$\Gamma_{t+j,t} \equiv \Gamma_{t+1}\,\Gamma_{t+2}\cdots\Gamma_{t+j}$
(with $\Gamma_{t,t}=1$).

\subsubsection{Value functions and optimality conditions}

Let $\widetilde{\mathcal{V}}_t$ denote the value (in utils) of a firm
that receives a favourable signal at date~$t$, and
$\mathcal{V}_t(P_{t-1}(\iota))$ the value of a firm that receives the
unfavourable signal. A firm receiving the unfavourable signal simply
follows the indexation rule~\eqref{eq:index_rule}, so its value at
date~$t$ depends on its inherited price $P_{t-1}(\iota)$. By contrast, a
firm receiving the favourable signal is purely forward-looking: since all
such firms share the same expectations, they all choose the same optimal
price~$\breve P_t$, and the value $\widetilde{\mathcal{V}}_t$ does not
depend on any endogenous state variable.

The Bellman equation for a firm receiving the favourable signal is:
\begin{equation}\label{eq:bellman_fav}
  \widetilde{\mathcal{V}}_t
  = \max_{\mathbf{P}}\left\{
    \Pi_t(\mathbf{P})
    + \beta\,\E_t\!\left[\frac{\Lambda_{t+1}}{\Lambda_t}
      \Bigl((1-\xi_p)\,\widetilde{\mathcal{V}}_{t+1}
      + \xi_p\,\mathcal{V}_{t+1}(\mathbf{P})\Bigr)\right]
  \right\},
\end{equation}
where $\Lambda_t$ is the nominal marginal utility of wealth (the
Lagrange multiplier on the household's budget constraint). The value of a
firm receiving the unfavourable signal is:
\begin{equation}\label{eq:bellman_unfav}
  \mathcal{V}_t\bigl(P_{t-1}(\iota)\bigr)
  = \Pi_t\!\left(\Gamma_t\,P_{t-1}(\iota)\right)
  + \beta\,\E_t\!\left[\frac{\Lambda_{t+1}}{\Lambda_t}
    \Bigl((1-\xi_p)\,\widetilde{\mathcal{V}}_{t+1}
    + \xi_p\,\mathcal{V}_{t+1}\bigl(\Gamma_t\,P_{t-1}(\iota)\bigr)
    \Bigr)\right].
\end{equation}

The first-order condition and the envelope condition for a firm receiving
the favourable signal give:
\begin{align}
  \Pi'_t(\breve P) + \beta\,\xi_p\,\E_t\!\left[
    \frac{\Lambda_{t+1}}{\Lambda_t}\,\mathcal{V}'_{t+1}(\breve P)\right]
  &= 0, \label{eq:foc_price}\\[4pt]
  \frac{\mathcal{V}'_t(P_{t-1}(\iota))}{\Gamma_t}
  = \Pi'_t\!\left(\Gamma_t\,P_{t-1}(\iota)\right)
  + \beta\,\xi_p\,\E_t\!\left[
    \frac{\Lambda_{t+1}}{\Lambda_t}\,
    \mathcal{V}'_{t+1}\bigl(\Gamma_t\,P_{t-1}(\iota)\bigr)\right].
    \label{eq:envelope_price}
\end{align}
The envelope condition~\eqref{eq:envelope_price} can be iterated forward.
Writing temporarily $\mathcal{P}$ for an inherited price, one step ahead:
\[
  \mathcal{V}'_{t+1}(\mathcal{P})
  = \Gamma_{t+1}\,\Pi'_{t+1}(\Gamma_{t+1}\,\mathcal{P})
  + \beta\,\xi_p\,\Gamma_{t+1}\,
    \E_{t+1}\!\left[\frac{\Lambda_{t+2}}{\Lambda_{t+1}}\,
    \mathcal{V}'_{t+2}(\Gamma_{t+1}\,\mathcal{P})\right].
\]
Iterating to infinity and taking the conditional expectation at~$t$:
\[
  \E_t\!\left[\mathcal{V}'_{t+1}(\mathcal{P})\right]
  = \E_t\!\left[\sum_{j=0}^{\infty}(\beta\,\xi_p)^j\,
    \Gamma_{t+1+j,t}\,\frac{\Lambda_{t+1+j}}{\Lambda_{t+1}}\,
    \Pi'_{t+1+j}\bigl(\Gamma_{t+1+j,t}\,\mathcal{P}\bigr)\right].
\]
Substituting back into the first-order
condition~\eqref{eq:foc_price} with $\mathcal{P} = \breve P_t$:
\begin{equation}\label{eq:optimal_price_sum}
  \E_t\!\left[\sum_{j=0}^{\infty}(\beta\,\xi_p)^j\,
  \Gamma_{t+j,t}\,\frac{\Lambda_{t+j}}{\Lambda_t}\,
  \Pi'_{t+j}\!\left(\Gamma_{t+j,t}\,\breve P_t\right)\right] = 0.
\end{equation}
The optimal price~$\breve P_t$ ensures that the sum (discounted and
probability-weighted) of the marginal profit over all future dates at
which the firm might still be using this price is zero. It also
guarantees the nullity of the expected marginal profits in the case
where the firm cannot re-optimise (by following the indexation
rule~\eqref{eq:index_rule}).

\subsubsection{Derivation of the $Z$ recursions}

To obtain a recursive formulation, we use the Kimball aggregator
multiplier $\Theta_t$, which satisfies~\eqref{eq:eq16}. The derivative
of the profit function with respect to~$\mathcal{P}$ is:
\[
  \Pi'_t(\mathcal{P}) =
  \frac{1-\theta_f(1+\psi_f)}{1+\psi_f}\,
  \left(\frac{\mathcal{P}}{P_t}\right)^{\!-(1+\psi_f)\theta_f}\!
  \Theta_t^{(1+\psi_f)\theta_f}\,Y_t
  + \theta_f\,
  \left(\frac{\mathcal{P}}{P_t}\right)^{\!-(1+\psi_f)\theta_f-1}\!
  \Theta_t^{(1+\psi_f)\theta_f}\,\nu^p_t\,mc_t\,Y_t
  + \frac{\psi_f}{1+\psi_f}\,Y_t.
\]
Substituting this expression into~\eqref{eq:optimal_price_sum} and
dividing through by $P_t^{*\,-(1+\psi_f)\theta_f}$, the infinite sum
decomposes into three terms that can each be written recursively. Denoting
by $\breve p_t = \breve P_t/P_t$ the optimal relative price:

\paragraph{Equation~15: Optimal relative price.}
\begin{equation}\label{eq:eq15}\tag{Eq.~15}
  \boxed{
  \breve p_t = \frac{\theta_f(1+\psi_f)}{\theta_f(1+\psi_f)-1}\,
  \frac{Z_{1,t}}{Z_{2,t}}
  + \frac{\psi_f}{\theta_f(1+\psi_f)-1}\,
  (\breve p_t)^{1+\theta_f(1+\psi_f)}\,
  \frac{Z_{3,t}}{Z_{2,t}}.}
\end{equation}
This is an \emph{implicit} equation for $\breve p_t$. The first term captures
the standard markup channel (marginal cost over revenue), while the
second term is the \textbf{Kimball correction}: it depends on
$(\breve p_t)^{1+\theta_f(1+\psi_f)}$ and vanishes when $\psi_f = 0$. In the
CES case ($\psi_f = 0$), the equation reduces to
$\breve p_t = [\theta_f/(\theta_f - 1)]\,Z_{1,t}/Z_{2,t}$, the familiar
constant-markup rule.

The three auxiliary variables $Z_{1,t}$, $Z_{2,t}$, $Z_{3,t}$ arise from
the three terms in~$\Pi'_t$ and can be represented recursively. Noting
that $\pi_{t+j}/P_t = \prod_{h=1}^{j}\pi_{t+h}$, the inflation factor
between $t$ and $t+j$ enters each recursion through the ratio
$\pi_{t+1}/\pi^{np}_t$ (i.e.\ the erosion of the firm's relative price
per period).

\paragraph{Equation~17: Recursion for $Z_{1,t}$ (marginal cost component).}
$Z_{1,t}$ captures the present-value effective marginal cost
$\nu^p_t\,mc_t$, weighted by the Kimball demand:
\begin{equation}\label{eq:eq17}\tag{Eq.~17}
  \boxed{
  Z_{1,t} = \lambda_t\,mc_t\,\nu^p_t\,(\Lambda^f_t)^{\theta_f(1+\psi_f)}\,Y_t
  + \beta\,\xi_p\,g_{t+1}^{1-\sigma_c}\,
    \left(\frac{\pi_{t+1}}{\pi^{np}_t}\right)^{\!\theta_f(1+\psi_f)}\!Z_{1,t+1}.}
\end{equation}
The term $(\Lambda^f_t)^{\theta_f(1+\psi_f)}$ arises from the Kimball
demand evaluated at the aggregate level. The discount factor
$(\pi_{t+1}/\pi^{np}_t)^{\theta_f(1+\psi_f)}$ reflects the erosion of
the firm's relative demand as its price drifts away from the optimum
(only indexed by $\pi^{np}_t$ while the aggregate price grows
by $\pi_{t+1}$).

\paragraph{Equation~18: Recursion for $Z_{2,t}$ (revenue component).}
$Z_{2,t}$ captures the present-value revenue, also weighted by Kimball
demand:
\begin{equation}\label{eq:eq18}\tag{Eq.~18}
  \boxed{
  Z_{2,t} = \lambda_t\,(\Lambda^f_t)^{\theta_f(1+\psi_f)}\,Y_t
  + \beta\,\xi_p\,g_{t+1}^{1-\sigma_c}\,
    \left(\frac{\pi_{t+1}}{\pi^{np}_t}\right)^{\!\theta_f(1+\psi_f)-1}\!Z_{2,t+1}.}
\end{equation}
The exponent $\theta_f(1+\psi_f)-1$ (rather than $\theta_f(1+\psi_f)$)
reflects that revenue is price$\,\times\,$quantity: the extra power
of~$-1$ comes from the price factor in the revenue integrand.

\paragraph{Equation~19: Recursion for $Z_{3,t}$ (Kimball correction).}
$Z_{3,t}$ is specific to the Kimball aggregator---it captures the
``linear'' part of the demand function (the $+\psi_f$ term
in~\eqref{eq:kimball_demand}):
\begin{equation}\label{eq:eq19}\tag{Eq.~19}
  \boxed{
  Z_{3,t} = \lambda_t\,Y_t
  + \beta\,\xi_p\,g_{t+1}^{1-\sigma_c}\,
    \frac{\pi^{np}_t}{\pi_{t+1}}\,Z_{3,t+1}.}
\end{equation}
Note the simpler discount factor: the ``linear'' component of the
Kimball demand contributes a term that only depends on the relative
price to the power~$1$ (not~$\theta_f(1+\psi_f)$), so the erosion
factor is just $\pi^{np}_t/\pi_{t+1}$.

\medskip
\noindent\textbf{Interpretation at steady state.}
At steady state, $Z_{1}/Z_{2} = mc = 1/\mu_f$ and $Z_{3}/Z_{2} = 1$,
so~\eqref{eq:eq15} reduces to:
\[
  1 = \frac{\theta_f(1+\psi_f)}{\theta_f(1+\psi_f)-1}\,\frac{1}{\mu_f}
  + \frac{\psi_f}{\theta_f(1+\psi_f)-1},
\]
which simplifies to $\mu_f = \theta_f/(\theta_f - 1)$, independently
of~$\psi_f$. The Kimball curvature does not affect the
\emph{steady-state} markup---it only affects its \emph{dynamics} by
making demand more elastic for firms that deviate from the average price.

\subsection{Linearised New Keynesian Phillips curve}
\label{sec:linearised_nkpc}

We now derive the linearised Phillips curve implied by the nonlinear
model, in order to clarify the role of the cost-push shock~$\nu^p_t$.
For transparency, we work in the CES limit
($\psi_f = 0$, $\Lambda^f_t = 1$), with zero trend inflation
($\bbar\pi = 1$) and a unit growth rate ($g = 1$); the Kimball
curvature modifies the slope of the Phillips curve but does not change
the qualitative conclusions. We denote log-deviations from steady
state by hats: $\hat x_t = \log(x_t/x^\ss)$.

\subsubsection{Log-linearisation of the price index}

In the CES case, the Kimball aggregator reduces to the Dixit--Stiglitz
price index:
\[
  1 = (1-\xi_p)\,(\breve p_t)^{1-\theta_f}
  + \xi_p\,\bigl(\pi^{np}_t/\pi_t\bigr)^{1-\theta_f}.
\]
Log-linearising around $\breve p^\ss = 1$ and ${\pi^{np}}^\ss/\pi^\ss = 1$:
\begin{equation}\label{eq:pstar_lin}
  \hat{\breve p}_t = \frac{\xi_p}{1-\xi_p}\,
  \bigl(\hat\pi_t - \gamma_p\,\hat\pi_{t-1}\bigr),
\end{equation}
using $\hat\pi^{np}_t = \gamma_p\,\hat\pi_{t-1}$ from the
indexation rule~\eqref{eq:eq12}.

\subsubsection{Log-linearisation of the $Z$ recursions}

With $\psi_f = 0$, the three recursions collapse to two (since $Z_3$
is absent). At steady state ($\nu^p = 1$),
$Z_1^\ss = \lambda^\ss\,mc^\ss\,Y^\ss/(1-\beta\xi_p)$ and
$Z_2^\ss = \lambda^\ss\,Y^\ss/(1-\beta\xi_p)$. Log-linearising
\eqref{eq:eq17} and \eqref{eq:eq18}:
\begin{align}
  \hat Z_{1,t} &= (1-\beta\xi_p)\,
    (\hat\lambda_t + \widehat{mc}_t + \hat\nu^p_t + \hat Y_t)
  + \beta\xi_p\,
    \bigl[\theta_f\,(\hat\pi_{t+1} - \gamma_p\hat\pi_t)
    + \hat Z_{1,t+1}\bigr],
    \label{eq:Z1lin}\\
  \hat Z_{2,t} &= (1-\beta\xi_p)\,
    (\hat\lambda_t + \hat Y_t)
  + \beta\xi_p\,
    \bigl[(\theta_f-1)(\hat\pi_{t+1} - \gamma_p\hat\pi_t)
    + \hat Z_{2,t+1}\bigr].
    \label{eq:Z2lin}
\end{align}
The log-linearised optimal price is
$\hat{\breve p}_t = \hat Z_{1,t} - \hat Z_{2,t}$ (the constant
$\log\mu_f$ cancels). Subtracting~\eqref{eq:Z2lin}
from~\eqref{eq:Z1lin}:
\begin{equation}\label{eq:pstar_foc}
  \hat{\breve p}_t = (1-\beta\xi_p)\,
    (\widehat{mc}_t + \hat\nu^p_t)
  + \beta\xi_p\,
    \bigl[\hat\pi_{t+1} - \gamma_p\hat\pi_t
    + \hat{\breve p}_{t+1}\bigr].
\end{equation}
Note that $\hat\lambda_t$ and $\hat Y_t$ cancel in the difference:
the reset price depends on the effective marginal
cost $\widehat{mc}_t + \hat\nu^p_t$, not on the level of demand.

\subsubsection{The hybrid New Keynesian Phillips curve}

Substituting~\eqref{eq:pstar_lin} into~\eqref{eq:pstar_foc} and
simplifying:
\begin{equation}\label{eq:nkpc}
  \boxed{
  \hat\pi_t = \frac{\gamma_p}{1+\beta\gamma_p}\,\hat\pi_{t-1}
  + \frac{\beta}{1+\beta\gamma_p}\,\E_t\hat\pi_{t+1}
  + \frac{\kappa}{1+\beta\gamma_p}\,
    \bigl(\widehat{mc}_t + \hat\nu^p_t\bigr),}
\end{equation}
where the slope of the Phillips curve is
\[
  \kappa \equiv \frac{(1-\xi_p)(1-\beta\xi_p)}{\xi_p}.
\]

\subsubsection{Comparison with the Smets--Wouters markup shock}
\label{sec:comparison_sw}

In the original SW model, which is formulated directly in linearised
form, the NKPC reads:
\[
  \hat\pi_t
  = \frac{\gamma_p}{1+\beta\gamma_p}\,\hat\pi_{t-1}
  + \frac{\beta}{1+\beta\gamma_p}\,\hat\pi_{t+1}
  + \frac{\kappa'}{1+\beta\gamma_p}\,\widehat{mc}_t
  + \eta^p_t,
\]
where $\eta^p_t$ is an exogenous ``price markup shock'' (an AR or
ARMA process) added directly to the Phillips curve, and $\kappa'$
may differ from $\kappa$ depending on the normalisation. Two
important differences with~\eqref{eq:nkpc} emerge:
\begin{enumerate}[nosep]
\item \textbf{Coefficient.}
  In~\eqref{eq:nkpc} the cost-push shock $\hat\nu^p_t$ enters with the
  \emph{same} slope coefficient $\kappa/(1+\beta\gamma_p)$ as the
  marginal cost.  This is because it only affects the
  \emph{contemporaneous} term in the $Z_1$ recursion and therefore
  picks up the weight $(1-\beta\xi_p)$ that the current period receives
  in the infinite discounted sum.

  By contrast, a ``true'' markup shock arising from a time-varying
  elasticity of substitution $\theta_{f,t}$ would enter through the
  \emph{front coefficient} $\mu_{f,t} = \theta_{f,t}/(\theta_{f,t}-1)$
  of the reset price equation (outside the infinite sum), yielding a
  different coefficient in the NKPC. The SW specification, where
  $\eta^p_t$ enters additively with its own normalisation, is not
  constrained by $\kappa$, and can therefore match a broader range of
  dynamics.

\item \textbf{Spillovers to other equations.}
  The SW markup shock $\eta^p_t$ only enters the Phillips curve.
  The cost-push shock $\nu^p_t$ enters only~$Z_1$ and therefore
  \emph{only} affects the Phillips curve as well---it is the
  closest structural proxy for the SW markup shock, free of
  side-effects on the resource constraint.
\end{enumerate}

\paragraph{Where the shocks appear.}
To confirm that the cost-push shocks act purely as Phillips-curve
shifters, we trace their appearances across all equilibrium
conditions:

\begin{center}
\renewcommand{\arraystretch}{1.2}
\begin{tabular}{lcc}
\hline
\textbf{Equation} & $\nu^p_t$ & $\nu^w_t$ \\
\hline
Eq.~\eqref{eq:eq17}: $Z_{1,t}$ (marginal cost, prices)
  & $\checkmark$ & --- \\
Eq.~\eqref{eq:eq28}: $H_{1,t}$ (household wage cost)
  & --- & $\checkmark$ \\
\hline
\end{tabular}
\end{center}

\noindent
The cost-push shocks $\nu^p_t$ and $\nu^w_t$ each enter a single
equation ($Z_1$ and $H_1$ respectively), making them pure Phillips-curve
shifters.  They appear nowhere else---not in the price/wage index
equations \eqref{eq:eq16}/\eqref{eq:eq27}, the dispersion indices, the
household's FOCs, the Taylor rule, the resource constraint, or the
efficient economy block.

In particular, the resource constraint~\eqref{eq:eq42} is
\textbf{undistorted}:
\[
  G_t + C_t + p^I_t\,I_t = Y_t,
\]
exactly as in the Smets--Wouters model, where the markup shocks enter
only the linearised Phillips curve (prices) and the wage equation
(wages). The cost-push shocks therefore carry no resource-constraint
side-effects.

\paragraph{Impulse responses.}
A positive innovation to $\nu^p_t$ (resp.\ $\nu^w_t$) raises inflation
(resp.\ wage inflation) through the Phillips-curve channel.  Since it
leaves the resource constraint undistorted, the output contraction
operates \emph{exclusively} through the interest-rate channel, exactly
as for the markup shocks in the linearised SW model.

\subsubsection{The Phillips-curve slope under Kimball aggregation}
\label{sec:kimball_slope}

With $\psi_f \neq 0$, the optimal price~\eqref{eq:eq15} is an
implicit equation involving $Z_3$.  The log-linearisation of
\eqref{eq:eq15} around $\breve p^\ss = 1$ gives
\[
  \underbrace{\frac{(\theta_f-1)}
    {\theta_f-1-\theta_f\psi_f}}_{\displaystyle a_1}
  \,(\hat Z_{1,t} - \hat Z_{2,t})
  \;+\;
  \underbrace{\frac{-\psi_f}
    {\theta_f-1-\theta_f\psi_f}}_{\displaystyle a_2}
  \,(\hat Z_{2,t} - \hat Z_{3,t})
  \;=\; \hat{\breve p}_t.
\]
The difference $\hat Z_{1,t} - \hat Z_{2,t}$ satisfies the same
recursion~\eqref{eq:pstar_foc} as in the CES case: the
$\Lambda^f_t{}^{\varepsilon_f}$ terms in $Z_1$ and $Z_2$ cancel in
the difference, and the cost-push shock $\hat\nu^p_t$ enters with
coefficient $(1-\beta\xi_p)$ as before.  The difference
$\hat Z_{2,t} - \hat Z_{3,t}$ involves only the Kimball multiplier
$\hat\Lambda^f_t$ ($\nu^p_t$ does not appear, since it only enters
$Z_1$).  Moreover, combining
the log-linearised price index~\eqref{eq:eq16} with the aggregate
price equation~\eqref{eq:eq21} shows that $\hat\Lambda^f_t = 0$
when starting from steady state, so
$\hat Z_{2,t} - \hat Z_{3,t}$ contributes only through inflation
expectations.

The resulting NKPC retains the hybrid form~\eqref{eq:nkpc} with a
modified slope
\begin{equation}\label{eq:kappa_kimball}
  \kappa^K
  = \frac{(1-\xi_p)(1-\beta\xi_p)}{\xi_p}
  \times \frac{\theta_f-1}{\theta_f-1-\theta_f\psi_f}
  = \frac{(1-\xi_p)(1-\beta\xi_p)}
    {\xi_p\,\bigl(1 + \varepsilon^K_f/(\theta_f-1)\bigr)},
\end{equation}
where $\varepsilon^K_f \equiv -\theta_f\psi_f > 0$ is the
effective Kimball curvature.  Since $\psi_f < 0$, the second factor
is less than one: the variable elasticity of demand creates
``strategic complementarity'' that dampens optimal price adjustments
and flattens the Phillips curve.  The cost-push shock still enters
with this same (flatter) slope, and the qualitative conclusions above
remain unchanged.  The same derivation
applies to the wage Phillips curve, with the analogous correction
$(\theta_s-1)/(\theta_s-1-\theta_s\psi_s)$.

\paragraph{Quantitative magnitude.}
With the Smets and Wouters (2007) posterior-mean calibration $\xi_p = 0.66$,
$\gamma_p = 0.24$, $\beta = 0.9975$ and the Kimball curvature
$\varepsilon^K_f = 10$, the
Kimball-corrected slope~\eqref{eq:kappa_kimball} evaluates to
\[
  \kappa^K \approx 0.060,
  \qquad
  \frac{\kappa^K}{1 + \beta\gamma_p}
  \approx 0.048.
\]
For wages, with $\xi_w = 0.70$, $\gamma_w = 0.58$ and the corresponding Kimball
curvature $\varepsilon^K_s = 10$:
\[
  \kappa^K_w \approx 0.043,
  \qquad
  \frac{\kappa^K_w}{1 + \beta\gamma_w}
  \approx 0.027.
\]
These values are consistent with the empirical literature on NKPC
estimation, which typically finds coefficients of $0.01$--$0.05$.

By contrast, the linearised markup shock in SW enters the NKPC with a
coefficient of~$1$ (or, equivalently, the shock is defined so that a
unit innovation raises inflation by its own magnitude).  Therefore,
\emph{for the cost-push shock $\nu^p_t$ to generate the same impact on
inflation as the SW markup shock, it would need to be roughly
$1/0.048 \approx 21$ times as volatile} (for prices) and
$1/0.027 \approx 37$ times as volatile (for wages).

\subsubsection{Structural interpretation and estimation}
\label{sec:mc_shock}

The cost-push shock~$\nu^p_t$ (and its wage-side
counterpart~$\nu^w_t$) has a natural structural interpretation: it
represents an exogenous increase in firms' input costs (e.g.\ imported
intermediates, energy prices, indirect taxes) that is orthogonal to
the quantity of goods produced.

\medskip\noindent
The cost-push formulation also avoids an identification tension. A
revenue-side markup proxy would enter two channels with very different
coefficients---the Phillips curve ($\approx 0.048$) and the resource
constraint ($\approx 1$)---forcing the estimator to reconcile the
inflationary signal with the output signal and potentially biasing
persistence or variance estimates.  The cost-push shock enters only the
Phillips curve, avoiding this tension by construction.

\subsection{Aggregate output and price distortion}

\paragraph{Equation~41: Aggregate production.}
Each intermediate firm~$\iota$ produces according to the Cobb--Douglas
technology:
\[
  y_t(\iota) = K^d_t(\iota){}^\alpha\,(A_t\,L^d_t(\iota))^{1-\alpha}.
\]
Since all firms face the same factor prices, they all choose the same
capital--labour ratio (see equation~\eqref{eq:eq11}):
\[
  \frac{K^d_t(\iota)}{L^d_t(\iota)} = \frac{K^d_t}{L^d_t}
  \qquad\text{for all }\iota\in[0,1].
\]
Thus we can write the output of firm~$\iota$ as:
\[
  y_t(\iota)
  = \left(\frac{K^d_t}{L^d_t}\right)^{\!\alpha}
    (A_t)^{1-\alpha}\,L^d_t(\iota)
  = \frac{K^d_t{}^\alpha\,(A_t\,L^d_t)^{1-\alpha}}{L^d_t}\,L^d_t(\iota).
\]
Denoting aggregate potential output $y_t \equiv K^d_t{}^\alpha
(A_t\,L^d_t)^{1-\alpha}$ (where $K^d_t = \int_0^1
K^d_t(\iota)\,\dd\iota$ and $L^d_t = \int_0^1
L^d_t(\iota)\,\dd\iota$), integrating over~$\iota$ and using the
Kimball demand function~\eqref{eq:kimball_demand}:
\[
  \int_0^1 y_t(\iota)\,\dd\iota
  = y_t\,\underbrace{\int_0^1
    \frac{L^d_t(\iota)}{L^d_t}\,\dd\iota}_{=1}
  = y_t.
\]
But because the aggregation technology is not linear ($G\neq$
identity), the sum of intermediate outputs differs from final output
$Y_t$. The relation is $\int_0^1 y_t(\iota)\,\dd\iota = \Delta^p_t
\cdot Y_t$ (integrating the demand function over~$\iota$), so:
\begin{equation}\label{eq:eq41}\tag{Eq.~41}
  \boxed{
  Y_t = \frac{K_t^{\alpha}\,(A_t\,L^d_t)^{1-\alpha}}{\Delta^p_t}.}
\end{equation}
The price dispersion term $\Delta^p_t \geq 1$ (with equality only at
the symmetric equilibrium) acts as a ``tax'' on aggregate production:
heterogeneous relative prices distort the allocation of inputs across
firms, reducing aggregate output below the level that would obtain if
all firms produced the same quantity. At the first order of
approximation (log-linearisation around the symmetric steady state
where $\Delta^p=1$), this term vanishes, but it matters at higher
orders.


%% ============================================================
\section{Labour market and wage setting}
\label{sec:wages}
%% ============================================================

The labour market mirrors the goods market structure exactly: a Kimball
aggregator combines differentiated labour varieties, and Calvo
wage-setting by unions replaces Calvo pricing by firms.  The derivation
proceeds in the same way as for the price-setting block
(Section~\ref{sec:prices}), replacing the final good producer by a
representative employment agency, intermediate goods firms by unions,
prices by wages, and the goods-market Kimball parameters
$(\theta_f,\psi_f)$ by their labour-market counterparts
$(\theta_s,\psi_s)$.  We therefore only state the resulting equations
and highlight the points where the wage block differs from the price
block.

\subsection{Employment agency: Kimball aggregation}

A representative \textbf{employment agency} aggregates differentiated
labour varieties $L_t(\varsigma)$ (supplied by unions $\varsigma\in[0,1]$) into a
homogeneous labour composite $L^s_t$ using a Kimball aggregator of the
same form as in the goods market:
\begin{equation}
  \int_0^1 G^w\!\left(\frac{L_t(\varsigma)}{L^s_t};\,\theta_s,\psi_s\right)\dd\varsigma = 1,
\end{equation}
where the function~$G^w$ has the same functional form as~$G$ in
equation~\eqref{eq:kimball_G}, with $\theta_s$ and~$\psi_s$ replacing
$\theta_f$ and~$\psi_f$.  The steady-state wage markup and the
curvature parameter are:
\begin{equation}
  \mu_s = \frac{\theta_s}{\theta_s - 1},
  \qquad\text{and}\qquad
  \psi_s = -\frac{\varepsilon_s}{\theta_s}.
\end{equation}

All aggregate conditions carry over from the price block by direct
analogy:

\paragraph{Equation~32: Aggregate wage normalisation.}
\begin{equation}\label{eq:eq32}\tag{Eq.~32}
  \boxed{
  \frac{\psi_s}{1+\psi_s}\,\tilde w_t + \frac{1}{1+\psi_s}\,\Lambda^s_t = 1.}
\end{equation}

\subsubsection{Wage dispersion}

The derivation mirrors the price dispersion case exactly
(Section~\ref{sec:prices}).

\paragraph{Definition of $\Delta^w_t$.}
The Kimball labour demand (the wage analogue
of~\eqref{eq:kimball_demand}) implies that the labour supplied by
union~$\varsigma$ relative to aggregate labour is:
\[
  \frac{l_t(\varsigma)}{L^s_t}
  = \frac{1}{1+\psi_s}
  \left[\left(\frac{W_t(\varsigma)/W_t}{\Upsilon_t}\right)^{\!-(1+\psi_s)\theta_s}
    + \psi_s\right].
\]
Integrating over all unions $\varsigma\in[0,1]$:
\begin{equation}\label{eq:deltaw_def}
  \Delta^w_t
  \equiv \int_0^1 \frac{l_t(\varsigma)}{L^s_t}\,\dd\varsigma
  = \frac{1}{1+\psi_s}
  \left[(\Lambda^s_t)^{\theta_s(1+\psi_s)}\!
    \int_0^1
    \left(\frac{W_t(\varsigma)}{W_t}\right)^{\!-\theta_s(1+\psi_s)}
    \!\dd\varsigma
    + \psi_s\right],
\end{equation}
where the integral $\int_0^1
(W_t(\varsigma)/W_t)^{-\theta_s(1+\psi_s)}\,\dd\varsigma$ is the
power mean of relative wages, denoted~$\nabla^w_t$.

\paragraph{Recursive expression for $\nabla^w_t$.}
Under the Calvo mechanism, the nominal wage $W_t(\varsigma)$ was set
optimally $j$ periods earlier with probability
$(1-\xi_w)\,\xi_w^j$. Since then it has been indexed by the cumulative
factor $\Omega_{t,t-j}$ (see~\eqref{eq:wage_index_rule}), while the
aggregate wage has grown by $W_t/W_{t-j}$. Thus:
\[
  \nabla^w_t
  = (1-\xi_w)\sum_{j=0}^{\infty}\xi_w^j
    \left(\frac{\Omega_{t,t-j}\,\breve W_{t-j}}{W_t}\right)^{\!-\theta_s(1+\psi_s)}.
\]
Splitting into $j=0$ (newly optimising unions, with relative wage
$\breve w_t$) and $j\geq 1$ (which gives $\nabla^w_{t-1}$ eroded by
one additional period of indexation versus actual wage growth):
\begin{equation}\label{eq:eq34}\tag{Eq.~34}
  \boxed{
  \nabla^w_t = (1-\xi_w)\,(\breve w_t)^{-\theta_s(1+\psi_s)}
  + \xi_w\,\nabla^w_{t-1}\,
    \left(\frac{\omega^{nw}_t}{\omega_t\,\pi_t}\right)^{\!-\theta_s(1+\psi_s)}.}
\end{equation}
The ratio $\omega^{nw}_t/(\omega_t\,\pi_t)$ captures the erosion of
each non-optimising union's relative wage: its nominal wage grows by
$\omega^{nw}_t$ (indexation, see~\eqref{eq:eq24}), while the aggregate
nominal wage grows by $\omega_t\,\pi_t$ (real wage growth times
inflation).

\paragraph{Equation~33: Wage dispersion index $\Delta^w_t$.}
Substituting the definition of $\nabla^w_t$ back
into~\eqref{eq:deltaw_def}:
\begin{equation}\label{eq:eq33}\tag{Eq.~33}
  \boxed{
  \Delta^w_t = \frac{\psi_s}{1+\psi_s}
  + \frac{1}{1+\psi_s}\,\nabla^w_t\,
    (\Lambda^s_t)^{\theta_s(1+\psi_s)}.}
\end{equation}
The structure is identical to the price dispersion
index~\eqref{eq:eq22}: the first term comes from the linear part of
the Kimball demand, and the second involves the power mean $\nabla^w_t$
scaled by the Kimball wage multiplier~$\Lambda^s_t$.

At steady state, $\breve w=1$, $\omega^{nw}=\omega\,\pi$, $\Lambda^s=1$,
so $\nabla^w = 1$ and $\Delta^w = 1$.

\paragraph{Equation~39: Household labour supply and employment agency supply.}
The wage dispersion creates a wedge between the household's total
labour~$L_t$ and the effective labour composite $L^s_t$ used by the
employment agency. Integrating the labour demand of
union~$\varsigma$ over all unions:
\[
  L_t
  = \int_0^1 l_t(\varsigma)\,\dd\varsigma
  = L^s_t\!\int_0^1 \frac{l_t(\varsigma)}{L^s_t}\,\dd\varsigma
  = \Delta^w_t\,L^s_t,
\]
so that:
\begin{equation}\label{eq:eq39}\tag{Eq.~39}
  \boxed{L^s_t = \frac{L_t}{\Delta^w_t}.}
\end{equation}
Since $\Delta^w_t \geq 1$, wage dispersion reduces the effective
labour available to firms relative to the total hours worked by
households, just as price dispersion reduces aggregate output relative
to aggregate inputs in~\eqref{eq:eq41}.

\subsection{Calvo wage setting}

In each period, a fraction $(1-\xi_w)$ of unions can re-optimise their
wage. The remaining fraction~$\xi_w$ indexes wages mechanically.

\paragraph{Equation~24: Wage indexation for non-optimising unions.}
Non-optimising unions adjust their nominal wage by:
\begin{equation}\label{eq:eq24}\tag{Eq.~24}
  \boxed{
  \omega^{nw}_t = \pi_t^{\gamma_w}\,\bbar\pi_t^{\,1-\gamma_w},}
\end{equation}
where $\gamma_w \in [0,1]$ is the degree of wage indexation to current
inflation.

\paragraph{Equation~25: Real wage growth.}
The real wage growth factor $\omega_t$ is defined as:
\begin{equation}\label{eq:eq25}\tag{Eq.~25}
  \boxed{\omega_t = \frac{w_t}{w_{t-1}}.}
\end{equation}

\paragraph{Equation~27: Aggregate Kimball multiplier (wages).}
\begin{equation}\label{eq:eq27}\tag{Eq.~27}
  \boxed{
  (\Lambda^s_t)^{1-\theta_s(1+\psi_s)}
  = (1-\xi_w)\,(\breve w_t)^{1-\theta_s(1+\psi_s)}
  + \xi_w\,
    \left(\frac{\omega^{nw}_t}{\omega_t\,\pi_t}\right)^{\!1-\theta_s(1+\psi_s)}
    \!(\Lambda^s_{t-1})^{1-\theta_s(1+\psi_s)}.}
\end{equation}

\paragraph{Equation~31: Average relative wage.}
\begin{equation}\label{eq:eq31}\tag{Eq.~31}
  \boxed{
  \tilde w_t = (1-\xi_w)\,\breve w_t
  + \xi_w\,\frac{\omega^{nw}_t}{\omega_t\,\pi_t}\,\tilde w_{t-1}.}
\end{equation}

\subsection{Optimal wage setting}

The derivation of the optimal wage follows the same value function
approach as for optimal pricing
(Section~\ref{sec:prices}), with unions replacing intermediate
firms, the nominal wage replacing the nominal price, and the employment
agency replacing the final good producer.

\subsubsection{The union's profit and the Calvo lottery}

Unions produce differentiated labour from the homogeneous labour supply
of households. The profit of a union offering nominal wage
$\mathcal{W}$ at date~$t$ is:
\[
  \mathscr{S}_t(\mathcal{W})
  = \bigl(\mathcal{W} - \nu^w_t\,W^m_t\bigr)\,
  \frac{1}{1+\psi_s}
  \left[\left(\frac{\mathcal{W}/W_t}{\Upsilon_t}\right)^{\!-(1+\psi_s)\theta_s}
    + \psi_s\right]\mathcal{L}_t,
\]
where $W^m_t$ is the nominal household wage, $W_t$ is the aggregate
nominal wage, $\Upsilon_t$ is the Kimball wage aggregator multiplier
(analogous to~$\Theta_t$ for prices), $\mathcal{L}_t$ is aggregate
labour demand, and $\nu^w_t$ is the wage cost-push shock
($\nu^w_t > 1$ raises the effective cost of labour for the union).

The cost-push shock~$\nu^w_t$ plays the same role on the wage side as
$\nu^p_t$ does on the price side (see Section~\ref{sec:prices}): it
enters only~$H_1$, making it a pure wage Phillips-curve shifter with no
effect on the resource constraint.

Each union is subject to the same Calvo lottery: with
probability~$\xi_w$ the union cannot re-optimise and indexes its nominal
wage by:
\begin{equation}\label{eq:wage_index_rule}
  W_t(\varsigma) = \Omega_t\,W_{t-1}(\varsigma),
  \qquad\text{where}\quad
  \Omega_t \equiv \frac{\mathcal{A}_{T,t}}{\mathcal{A}_{T,t-1}}\,
    \pi_{t-1}^{\gamma_w}\,\bbar\pi^{\,1-\gamma_w},
\end{equation}
i.e.\ a combination of lagged inflation, the inflation target, and the
growth rate of neutral technology $\mathcal{A}_{T,t}$. More generally,
$\Omega_{t+j,t} \equiv \Omega_{t+1}\cdots\Omega_{t+j}$ (with
$\Omega_{t,t}=1$).

\subsubsection{Value functions and optimality conditions}

Let $\widetilde{\mathcal{U}}_t$ denote the value of a union that
receives a favourable signal at date~$t$, and
$\mathcal{U}_t(W_{t-1}(\varsigma))$ the value of a union receiving the
unfavourable signal. As in the price-setting case, the value of an
optimising union is purely forward-looking, while the value of a
non-optimising union depends on its inherited wage.

The Bellman equation for a union receiving the favourable signal is:
\begin{equation}\label{eq:bellman_wage_fav}
  \widetilde{\mathcal{U}}_t
  = \max_{\mathbf{W}}\left\{
    \mathscr{S}_t(\mathbf{W})
    + \beta\,\E_t\!\left[\frac{\Lambda_{t+1}}{\Lambda_t}
      \Bigl((1-\xi_w)\,\widetilde{\mathcal{U}}_{t+1}
      + \xi_w\,\mathcal{U}_{t+1}(\mathbf{W})\Bigr)\right]
  \right\},
\end{equation}
and the value of a union receiving the unfavourable signal is:
\begin{equation}\label{eq:bellman_wage_unfav}
  \mathcal{U}_t\bigl(W_{t-1}(\varsigma)\bigr)
  = \mathscr{S}_t\!\left(\Omega_t\,W_{t-1}(\varsigma)\right)
  + \beta\,\E_t\!\left[\frac{\Lambda_{t+1}}{\Lambda_t}
    \Bigl((1-\xi_w)\,\widetilde{\mathcal{U}}_{t+1}
    + \xi_w\,\mathcal{U}_{t+1}\bigl(\Omega_t\,W_{t-1}(\varsigma)\bigr)
    \Bigr)\right].
\end{equation}

The first-order condition and the envelope condition give:
\begin{align}
  \mathscr{S}'_t(\breve W) + \beta\,\xi_w\,\E_t\!\left[
    \frac{\Lambda_{t+1}}{\Lambda_t}\,\mathcal{U}'_{t+1}(\breve W)\right]
  &= 0, \label{eq:foc_wage}\\[4pt]
  \frac{\mathcal{U}'_t(W_{t-1}(\varsigma))}{\Omega_t}
  = \mathscr{S}'_t\!\left(\Omega_t\,W_{t-1}(\varsigma)\right)
  + \beta\,\xi_w\,\E_t\!\left[
    \frac{\Lambda_{t+1}}{\Lambda_t}\,
    \mathcal{U}'_{t+1}\bigl(\Omega_t\,W_{t-1}(\varsigma)\bigr)\right].
    \label{eq:envelope_wage}
\end{align}
Iterating the envelope condition~\eqref{eq:envelope_wage} forward and
substituting into~\eqref{eq:foc_wage}:
\begin{equation}\label{eq:optimal_wage_sum}
  \E_t\!\left[\sum_{j=0}^{\infty}(\beta\,\xi_w)^j\,
  \Omega_{t+j,t}\,\frac{\Lambda_{t+j}}{\Lambda_t}\,
  \mathscr{S}'_{t+j}\!\left(\Omega_{t+j,t}\,\breve W_t\right)\right] = 0.
\end{equation}
The optimal nominal wage~$\breve W_t$ ensures the nullity of the sum
(discounted and probability-weighted) of the marginal union profits over
all future dates at which the union might still be using this wage.

\subsubsection{Derivation of the $H$ recursions}

The derivative of the union's profit with respect to~$\mathcal{W}$ has
the same three-term structure as in the price case: the cost-push
shock~$\nu^w_t$ multiplies the household wage in the cost term (yielding
$H_{1,t}$), while the revenue terms ($H_{2,t}$ and $H_{3,t}$) carry no
shock. Substituting
into~\eqref{eq:optimal_wage_sum} and dividing through by
$W_t^{*\,-(1+\psi_s)\theta_s}$, we obtain the optimal relative
wage $\breve w_t = \breve W_t/W_t$ as an implicit function of three recursion
variables $H_{1,t}$, $H_{2,t}$, $H_{3,t}$.

The household real wage $w^h_t$ that a member of union~$\varsigma$ receives is
related to the gross wage through the Kimball markup: at steady state,
$w^h = w/\mu_s$.

\paragraph{Equation~26: Optimal relative wage.}
\begin{equation}\label{eq:eq26}\tag{Eq.~26}
  \boxed{
  \breve w_t = \frac{\theta_s(1+\psi_s)}{\theta_s(1+\psi_s)-1}\,
  \frac{H_{1,t}}{H_{2,t}}
  + \frac{\psi_s}{\theta_s(1+\psi_s)-1}\,
  (\breve w_t)^{1+\theta_s(1+\psi_s)}\,
  \frac{H_{3,t}}{H_{2,t}}.}
\end{equation}
This parallels the optimal price equation~\eqref{eq:eq15}. The ratio
$H_{1,t}/H_{2,t}$ captures the desired wage markup (household wage over
gross wage), and $H_{3,t}$ corrects for Kimball curvature. In the CES
case ($\psi_s = 0$), the equation reduces to
$\breve w_t = [\theta_s/(\theta_s-1)]\,H_{1,t}/H_{2,t}$.

\paragraph{Equation~28: Recursion for $H_{1,t}$.}
$H_{1,t}$ captures the present-value effective household wage cost
$\nu^w_t\,w^h_t$, weighted by Kimball labour demand:
\begin{equation}\label{eq:eq28}\tag{Eq.~28}
  \boxed{
  H_{1,t} = \lambda_t\,w^h_t\,\nu^w_t\,(\Lambda^s_t)^{\theta_s(1+\psi_s)}\,L^s_t
  + \beta\,\xi_w\,g_{t+1}^{1-\sigma_c}\,
    \left(\frac{\omega_{t+1}\,\pi_{t+1}}{\omega^{nw}_t}\right)^{\!\theta_s(1+\psi_s)}
    \!H_{1,t+1}.}
\end{equation}

\paragraph{Equation~29: Recursion for $H_{2,t}$.}
$H_{2,t}$ captures the present-value revenue of the union, weighted
by Kimball labour demand:
\begin{equation}\label{eq:eq29}\tag{Eq.~29}
  \boxed{
  H_{2,t} = \lambda_t\,w_t\,
    (\Lambda^s_t)^{\theta_s(1+\psi_s)}\,L^s_t
  + \beta\,\xi_w\,g_{t+1}^{1-\sigma_c}\,
    \left(\frac{\omega_{t+1}\,\pi_{t+1}}{\omega^{nw}_t}\right)^{\!\theta_s(1+\psi_s)-1}
    \!H_{2,t+1}.}
\end{equation}
The exponent $\theta_s(1+\psi_s)-1$ (rather than $\theta_s(1+\psi_s)$)
has the same origin as in the price block: the extra power of~$-1$ comes
from the wage factor in the revenue integrand.

\paragraph{Equation~30: Recursion for $H_{3,t}$.}
$H_{3,t}$ is the Kimball correction, capturing the linear part of the
labour demand function:
\begin{equation}\label{eq:eq30}\tag{Eq.~30}
  \boxed{
  H_{3,t} = \lambda_t\,w_t\,L^s_t
  + \beta\,\xi_w\,g_{t+1}^{1-\sigma_c}\,
    \frac{\omega^{nw}_t}{\pi_{t+1}\,\omega_{t+1}}\,H_{3,t+1}.}
\end{equation}
As in the price block, the simpler discount factor reflects that the
linear component of the Kimball demand only depends on the relative wage
to the power~$1$.

\medskip
\noindent\textbf{Interpretation at steady state.}
At steady state, $H_1/H_2 = w^h/w = 1/\mu_s$ and $H_3/H_2 = 1$, so
equation~\eqref{eq:eq26} gives:
\[
  1 = \frac{\theta_s(1+\psi_s)}{\theta_s(1+\psi_s)-1}\,\frac{1}{\mu_s}
  + \frac{\psi_s}{\theta_s(1+\psi_s)-1},
\]
consistent with $\mu_s = \theta_s/(\theta_s - 1)$.


%% ============================================================
\section{Monetary policy}
\label{sec:monetary}
%% ============================================================

\paragraph{Equation~35: Taylor rule.}
The central bank sets the nominal interest rate according to a generalised
Taylor rule with interest rate smoothing:
\begin{equation}\label{eq:eq35}\tag{Eq.~35}
  \boxed{
  R_t = R_{t-1}^{\rho_R}
  \left[
  \bbar R\,
  \left(\frac{\pi_{t-1}}{\bbar\pi_t}\right)^{\!r_\pi}\!
  \mathcal{G}_t^{\,r_y}\,
  \left(\frac{\pi_t/\bbar\pi_t}{\pi_{t-1}/\bbar\pi_{t-1}}\right)^{\!r_{\Delta\pi}}\!
  \left(\frac{\mathcal{G}_t}{\mathcal{G}_{t-1}}\right)^{\!r_{\Delta y}}
  \right]^{1-\rho_R}
  \!\!\varepsilon^R_t,}
\end{equation}
where:
\begin{itemize}[nosep]
  \item $\rho_R \in [0,1)$ is the interest rate smoothing parameter;
  \item $\bbar R$ is the steady-state nominal interest rate;
  \item $r_\pi > 1$ is the response to lagged inflation deviations from target;
  \item $r_y > 0$ is the response to the output gap~$\mathcal{G}_t$;
  \item $r_{\Delta\pi}$ is the response to inflation growth;
  \item $r_{\Delta y}$ is the response to output gap growth;
  \item $\varepsilon^R_t$ is the monetary policy shock.
\end{itemize}


%% ============================================================
\section{Government}
\label{sec:government}
%% ============================================================

\paragraph{Equation~36: Output gap.}
The output gap is defined as the ratio of actual GDP to efficient GDP
(the level that would prevail under flexible prices and wages):
\begin{equation}\label{eq:eq36}\tag{Eq.~36}
  \boxed{\mathcal{G}_t = \frac{Y_t}{Y^e_t},}
\end{equation}
where $Y^e_t$ denotes efficient output (see \Cref{sec:efficient}).

\paragraph{Equation~37: Public spending.}
Government expenditure is a time-varying share of GDP:
\begin{equation}\label{eq:eq37}\tag{Eq.~37}
  \boxed{G_t = \varepsilon^g_t\,Y_t.}
\end{equation}


%% ============================================================
\section{Market clearing}
\label{sec:clearing}
%% ============================================================

\paragraph{Equation~38: Capital market.}
Effective capital (services) equals the utilised physical stock, detrended
by the growth rate:
\begin{equation}\label{eq:eq38}\tag{Eq.~38}
  \boxed{K_t = z_t\,\frac{\bbar K_{t-1}}{g_t}.}
\end{equation}

\paragraph{Equation~40: Labour market.}
The employment agency's labour supply equals labour demand from
intermediate firms:
\begin{equation}\label{eq:eq40}\tag{Eq.~40}
  \boxed{L^s_t = L^d_t.}
\end{equation}

\paragraph{Equation~42: Goods market.}
The aggregate resource constraint equates final demand to output:
\begin{equation}\label{eq:eq42}\tag{Eq.~42}
  \boxed{
  G_t + C_t + p^I_t\,I_t = Y_t.}
\end{equation}
Government spending, consumption and investment (valued at the relative
price of investment~$p^I_t$) exhaust output. Because the cost-push
shocks enter only the Phillips and wage curves, no markup wedge appears
here, and the constraint takes the same undistorted form as in the
linearised Smets--Wouters model.


%% ============================================================
\section{Exogenous processes (Equations~43--59)}
\label{sec:shocks}
%% ============================================================

By default, each of the 17 exogenous driving forces follows an AR(1)
process in levels (Section~\ref{sec:arma} generalises this to ARMA$(p,q)$).
Let $X_t$ denote a generic shock with steady-state value
$\bbar X$, autocorrelation parameter $\rho_X$, and innovation
$\epsilon_{X,t}$. The law of motion is:%
\footnote{For the growth-rate shock $g_t$
(\texttt{ProductionEfficiencyGrowth}), the normalisation uses
$(1+\bbar g)$ instead of $\bbar g$.}
\begin{equation}\label{eq:shock_general}\tag{Eq.~43--59}
  \boxed{
  \frac{X_t}{\bbar X}
  = \left(\frac{X_{t-1}}{\bbar X}\right)^{\!\rho_X}
  \exp(\sigma_X\,\epsilon_{X,t}),
  \qquad \epsilon_{X,t} \sim \mathcal{N}(0,1).}
\end{equation}
In log-deviations $\hat x_t = \ln(X_t/\bbar X)$, this is:
\[
  \hat x_t = \rho_X\,\hat x_{t-1} + \sigma_X\,\epsilon_{X,t},
\]
a standard AR(1) process.

\medskip
\noindent The 17 shocks are listed in \Cref{tab:shocks}, with their
Smets--Wouters target persistence $\rho_{SW}$, unconditional standard
deviation $\sqrt{V_{SW}}$, and ARMA order $(p,q)$. For an AR(1) the innovation
standard deviation is $\sigma_\varepsilon = \sqrt{V_{SW}}\,\sqrt{1-\rho_{SW}^2}$.
Shocks with $\sqrt{V_{SW}} = 0$ are kept in the model but switched off.

\begin{table}[h]
\centering\small
\begin{tabular}{llccc}
\toprule
\textbf{Shock} & \textbf{Symbol} & $\rho_{SW}$ & $\sqrt{V_{SW}}$ & $(p,q)$ \\
\midrule
Technology growth    & $g_t$              & $0.39$ & $0$      & $(1,0)$ \\
Technology level     & $A_t$              & $0.95$ & $0.0045$ & $(1,0)$ \\
Consumption tax      & $\tau^C_t$         & $0.00$ & $0$      & $(1,0)$ \\
Labour income tax    & $\tau^W_t$         & $0.00$ & $0$      & $(1,0)$ \\
Income tax           & $\tau^R_t$         & $0.00$ & $0$      & $(1,0)$ \\
Risk premium         & $\varepsilon^B_t$  & $0.22$ & $0.0024$ & $(1,0)$ \\
Labour supply        & $\varepsilon^L_t$  & $0.94$ & $0$      & $(1,0)$ \\
Inv.\ relative price & $p^I_t$           & $0.71$ & $0$      & $(1,0)$ \\
Inv.\ efficiency     & $\varepsilon^I_t$  & $0.71$ & $0.0045$ & $(1,0)$ \\
Employer labour tax  & $\tau^L_t$         & $0.00$ & $0$      & $(1,0)$ \\
Price cost-push      & $\nu^p_t$          & $0.90$ & $0.0014$ & $(1,1)$ \\
Wage cost-push       & $\nu^w_t$          & $0.97$ & $0.0024$ & $(1,1)$ \\
Inflation target     & $\bbar\pi_t$       & $0.00$ & $0$      & $(1,0)$ \\
Monetary policy      & $\varepsilon^R_t$  & $0.15$ & $0.0024$ & $(1,0)$ \\
Public spending      & $\varepsilon^g_t$  & $0.97$ & $0.0053$ & $(1,0)$ \\
Habit                & $\eta^s_t$         & $0.98$ & $0$      & $(1,0)$ \\
Preference           & $\varsigma_t$      & $0.90$ & $0$      & $(1,0)$ \\
\bottomrule
\end{tabular}
\caption{Exogenous shock processes: Smets--Wouters target persistence and
unconditional standard deviation, and ARMA order. Shocks with $\sqrt{V_{SW}}=0$
are kept in the model but switched off; the two mark-ups are ARMA(1,1), all
other active shocks AR(1).}\label{tab:shocks}
\end{table}

\subsection{Generalisation to ARMA$(p,q)$}
\label{sec:arma}

The AR(1) law~\eqref{eq:shock_general} is only the default. Because the model
is assembled programmatically with \texttt{modBuilder}, any shock can instead be
declared as an ARMA$(p,q)$ process, the orders $p$ and $q$ being chosen per
shock. In log deviations $\hat x_t = \ln(X_t/\bbar X)$,
\begin{equation}\label{eq:arma}
  \hat x_t = \sum_{i=1}^{p}\phi_i\,\hat x_{t-i}
           + \sigma_\varepsilon\Big(\varepsilon_t
             + \sum_{j=1}^{q}\theta_j\,\varepsilon_{t-j}\Big),
  \qquad \varepsilon_t\sim\mathcal{N}(0,1),
\end{equation}
written in the same multiplicative form as~\eqref{eq:shock_general},
\[
  \frac{X_t}{\bbar X}
  = \prod_{i=1}^{p}\Big(\frac{X_{t-i}}{\bbar X}\Big)^{\!\phi_i}
    \exp\!\Big(\sigma_\varepsilon\big(\varepsilon_t
      + \textstyle\sum_{j=1}^{q}\theta_j\,\varepsilon_{t-j}\big)\Big).
\]
The lagged innovations $\varepsilon_{t-j}$ are written directly: Dynare's
preprocessor introduces the corresponding \texttt{AUX\_EXO\_LAG} auxiliary
variables, so no manual bookkeeping is required.

\paragraph{Calibration by moment matching.}
The coefficients are chosen so that the ARMA process reproduces \emph{two}
moments of the corresponding Smets--Wouters shock: its unconditional variance
$\gamma(0)=V_{SW}$ and its first-order autocorrelation $\rho(1)=\rho_{SW}$ (for
an AR(1), $V_{SW}=\sigma_{SW}^2/(1-\rho_{SW}^2)$ and $\rho_{SW}$ is the
persistence). Two facts make this well posed:
\begin{itemize}[nosep]
\item the unconditional variance is a pure scale, so $\sigma_\varepsilon^2$ is
  set last to hit $V_{SW}$ exactly;
\item $\rho(1)$ is scale-free, so it is the only constraint that binds the AR/MA
  coefficients.
\end{itemize}
The remaining coefficients are therefore free, and that freedom is the point:
for $p+q\ge 2$ the process matches the variance and first-order persistence of
the SW shock while having a \emph{different} higher-order autocorrelation
structure (hence a different spectrum and different impulse responses).

\paragraph{Random generation.}
The free coefficients are drawn at random. We sample $p+q-1$ characteristic
roots with moduli in a configurable interval $[a,b]\subset(0,1)$ --- which
guarantees stationarity of the AR part and invertibility of the MA part --- and
solve the last root so that $\rho(1)=\rho_{SW}$; $\sigma_\varepsilon^2$ then
matches the variance. Not every target is attainable at a given order: an
ARMA(1,1) satisfies $\rho(1)\le(1+\phi)/2$, so reaching the wage mark-up
persistence $\rho_{SW}=0.97$ requires roots close to the unit circle ($b$ near
$1$); the generator reports infeasibility otherwise.

\paragraph{Autocovariances.}
Matching relies on the autocovariances of the ARMA$(p,q)$, obtained from the
MA$(\infty)$ representation $\hat x_t=\sum_{j\ge0}\psi_j\,\varepsilon_{t-j}$ with
$\psi_0=1$ and $\psi_j=\theta_j+\sum_{i}\phi_i\,\psi_{j-i}$, as
$\gamma(k)=\sigma_\varepsilon^2\sum_{j\ge0}\psi_j\,\psi_{j+k}$.

\paragraph{Implementation.}
Three routines support this (in the \texttt{matlab/} subdirectory, called from
\texttt{sw.m}): \texttt{arma\_autocov} (autocovariances), \texttt{arma\_random}
(the random moment-matching generator) and \texttt{declare\_arma\_shock} (which
writes the equation and declares the parameters $\phi_i,\theta_j,\sigma_\varepsilon$
and the innovation through the \texttt{modBuilder} API); they are checked in
\texttt{validate\_arma.m}. Setting
$(p,q)=(1,0)$ recovers the AR(1) baseline exactly. In the calibration shipped
here only the two mark-up shocks are declared as ARMA(1,1); for the price
mark-up this yields $\phi_1=0.903,\ \theta_1=-0.013$ (so that $\rho(1)=0.90$),
and for the wage mark-up $\phi_1=0.959,\ \theta_1=0.200$ ($\rho(1)=0.97$).


%% ============================================================
\section{The efficient economy (Equations~60--78)}
\label{sec:efficient}
%% ============================================================

The \textbf{efficient economy} is a version of the model in which nominal
rigidities (Calvo pricing and wage setting) are removed, but all real
frictions are kept:
\begin{itemize}[nosep]
  \item External habit formation with stochastic habit;
  \item Investment adjustment costs;
  \item Variable capacity utilisation and endogenous depreciation;
  \item All structural shocks.
\end{itemize}
The efficient economy is used to compute the output gap
$\mathcal{G}_t = Y_t/Y^e_t$ that enters the Taylor rule.

\subsection{Household equations (Equations~60--68)}

The household first-order conditions are identical to those in the sticky-price
economy (\Cref{sec:households}), with two modifications:
\begin{enumerate}[nosep]
  \item The nominal interest rate $R_t/\pi_{t+1}$ is replaced by the real
    interest rate $r^e_t$ (\texttt{EfficientRealInterestFactor}).
  \item All efficient variables carry a superscript~$e$.
\end{enumerate}

\paragraph{Equation~60: Efficient consumption FOC.}
\begin{equation}\tag{Eq.~60}
  \bigl(C^e_t - \eta\,\eta^s_t\,C^e_{t-1}/g_t + \ln\varsigma_t\bigr)^{-\sigma_c}
  \exp\!\Bigl(\tfrac{(\sigma_c-1)\varepsilon^L_t}{1+\sigma_l}
  (L^e_t/\bbar L)^{1+\sigma_l}\Bigr)
  = \tau^C_t\,\lambda^e_t.
\end{equation}

\paragraph{Equation~61: Efficient Euler equation.}
\begin{equation}\tag{Eq.~61}
  \lambda^e_t\,\tau^R_t
  = \beta\,\varepsilon^B_t\,r^e_t\,g_{t+1}^{-\sigma_c}\,
    \lambda^e_{t+1}\,\tau^R_{t+1}.
\end{equation}

\paragraph{Equations~62--68.} The efficient labour FOC, investment FOC,
depreciation rate, marginal depreciation, capacity utilisation FOC, capital
FOC, and capital accumulation are all identical in structure to
Equations~3--9, with all variables replaced by their efficient
counterparts.

\subsection{Price and wage decisions (Equations~69--72)}

Without nominal rigidities, all firms and unions re-optimise every period.
The Kimball aggregator still applies, so steady-state markups remain at
$\mu_f$ and $\mu_s$.

\paragraph{Equation~69: Efficient factor price frontier.}
\begin{equation}\tag{Eq.~69}
  \frac{\tau^L_t\,w^e_t\,L^{e,d}_t}{r^{e,k}_t\,K^e_t}
  = \frac{1-\alpha}{\alpha}.
\end{equation}

\paragraph{Equation~70: Efficient marginal cost.}
\begin{equation}\tag{Eq.~70}
  mc^e_t = A_t^{\,\alpha-1}\,
  \Bigl(\frac{r^{e,k}_t}{\alpha}\Bigr)^{\!\alpha}\,
  \Bigl(\frac{\tau^L_t\,w^e_t}{1-\alpha}\Bigr)^{\!1-\alpha}.
\end{equation}

\paragraph{Equation~71: Efficient price decision.}
With all firms re-optimising every period ($\xi_p = 0$ effectively), and all
setting the same price, the Kimball optimal price condition reduces to:
\begin{equation}\tag{Eq.~71}
  \frac{\theta_f(1+\psi_f)}{\theta_f(1+\psi_f)-1}\,mc^e_t
  + \frac{\psi_f}{\theta_f(1+\psi_f)-1}
  = 1.
\end{equation}
This pins down the efficient marginal cost at its steady-state value
$mc^e = 1/\mu_f$.

\paragraph{Equation~72: Efficient wage decision.}
Similarly, with all unions re-optimising:
\begin{equation}\tag{Eq.~72}
  \frac{\theta_s(1+\psi_s)}{\theta_s(1+\psi_s)-1}\,
  \frac{w^{e,h}_t}{w^e_t}
  + \frac{\psi_s}{\theta_s(1+\psi_s)-1}
  = 1,
\end{equation}
which gives $w^{e,h}/w^e = 1/\mu_s$.

\subsection{Market clearing (Equations~73--78)}

\paragraph{Equation~73: Efficient labour market.}
Without wage dispersion ($\Delta^w = 1$), labour supply equals demand directly:
\begin{equation}\tag{Eq.~73}
  L^{e,d}_t = L^e_t.
\end{equation}

\paragraph{Equation~74: Efficient capital market.}
\begin{equation}\tag{Eq.~74}
  K^e_t = z^e_t\,\frac{\bbar K^e_{t-1}}{g_t}.
\end{equation}

\paragraph{Equation~75: Efficient goods market.}
Without markup shocks distorting the resource constraint:
\begin{equation}\tag{Eq.~75}
  C^e_t + p^I_t\,I^e_t = (1 - \varepsilon^g_t)\,Y^e_t.
\end{equation}

\paragraph{Equation~76: Efficient aggregate production.}
Without price dispersion ($\Delta^p = 1$):
\begin{equation}\tag{Eq.~76}
  Y^e_t = (K^e_t)^{\alpha}\,(A_t\,L^{e,d}_t)^{1-\alpha}.
\end{equation}

\paragraph{Equations~77--78: Efficient investment costs.}
These are identical in form to Equations~10--11, with all variables
replaced by their efficient counterparts.


%% ============================================================
\section{Steady state}
\label{sec:steady_state}
%% ============================================================

The steady state is computed by setting all shocks to their long-run
values and all growth rates to their balanced-growth values. Key
steady-state relationships include:

\begin{enumerate}[nosep]
  \item The \textbf{discount factor} is pinned down by the Euler
    equation~\eqref{eq:eq2}:
    \[
      \beta = \frac{\bbar\pi}{\bbar\varepsilon^B\,\bbar R}\,
      (1+\bbar g)^{\sigma_c}.
    \]
  \item The \textbf{steady-state capital rental rate} from the capital
    FOC~\eqref{eq:eq8}:
    \[
      \bbar r^k = \frac{(1+\bbar g)^{\sigma_c} - \beta(1-\bbar\delta)}
      {\beta\,\bbar\tau^R\,\bbar z}.
    \]
  \item The \textbf{capital share} from the production function:
    \[
      \alpha = 1 - \bbar\tau^L\,\mu_f\,\frac{\bbar w\,\bbar L}{\bbar Y}.
    \]
  \item The \textbf{steady-state gross wage}:
    \[
      \bbar w = \frac{(1-\alpha)}{\bbar\tau^L}\,
      \Bigl(\frac{1}{\mu_f}\,
      \Bigl(\frac{\alpha}{\bbar r^k}\Bigr)^{\!\alpha}\Bigr)^{\!1/(1-\alpha)}\,
      \bbar A.
    \]
  \item The \textbf{investment-to-GDP ratio}:
    \[
      \frac{\bbar I}{\bbar Y}
      = \frac{\bbar g + \bbar\delta}{\bbar z}\,
        \frac{(\bbar K/\bbar L)^{1-\alpha}}{\bbar A^{1-\alpha}},
    \]
    where $\bbar K/\bbar L = \bbar\tau^L\,\bbar w\,\alpha/[(1-\alpha)\bbar r^k]$.
  \item The \textbf{consumption-to-GDP ratio}:
    \[
      \frac{\bbar C}{\bbar Y}
      = 1 - \bbar\varepsilon^g - \frac{\bbar I}{\bbar Y}.
    \]
\end{enumerate}

At steady state, all relative prices equal unity ($\breve p = \tilde p =
\Lambda^f = 1$), all dispersion indices equal unity ($\Delta^p = \Delta^w
= \nabla^p = \nabla^w = 1$), Tobin's $Q = 1$, and the output gap
$\mathcal{G} = 1$.


%% ============================================================
\section{Summary of the equation system}
\label{sec:summary}
%% ============================================================

The full model consists of:
\begin{itemize}[nosep]
  \item \textbf{11 household equations} (Eq.~1--11): consumption, Euler,
    labour, investment, depreciation (2), capacity utilisation, capital
    (Tobin's~$Q$), capital accumulation, investment cost (2);
  \item \textbf{12 price-setting equations} (Eq.~12--23): indexation,
    factor frontier, marginal cost, optimal price, Kimball multiplier,
    $Z_1$, $Z_2$, $Z_3$ recursions, average price, aggregate price,
    price dispersion, $\nabla^p$ recursion;
  \item \textbf{11 wage-setting equations} (Eq.~24--34): indexation, real
    wage growth, optimal wage, Kimball multiplier, $H_1$, $H_2$, $H_3$
    recursions, average wage, aggregate wage, wage dispersion, $\nabla^w$
    recursion;
  \item \textbf{3 policy/government equations} (Eq.~35--37): Taylor rule,
    output gap, public spending;
  \item \textbf{5 market clearing equations} (Eq.~38--42): capital market,
    labour supply--agency, labour demand, aggregate production, goods
    market;
  \item \textbf{17 exogenous processes} (Eq.~43--59): AR(1) laws of
    motion for all structural shocks;
  \item \textbf{19 efficient-economy equations} (Eq.~60--78): the
    flexible-price/wage counterpart.
\end{itemize}

This gives a total of 78 equations for 78 endogenous variables (plus the
17~exogenous processes), forming a complete nonlinear rational expectations
system that can be solved using perturbation methods in Dynare.

\end{document}
